% appendix-info-theory.tex
% Standalone appendix — include from retrospective.tex via:
%   \appendix
%   % appendix-info-theory.tex
% Standalone appendix — include from retrospective.tex via:
%   \appendix
%   % appendix-info-theory.tex
% Standalone appendix — include from retrospective.tex via:
%   \appendix
%   % appendix-info-theory.tex
% Standalone appendix — include from retrospective.tex via:
%   \appendix
%   \input{appendix-info-theory}
%
% Requires in the preamble: amsmath, amssymb, tabularx, booktabs

\section{Information-Theoretic View}
\label{app:info-theory}

This appendix applies information-theoretic measures to the real-time
web communication stack described in this document. We quantify channel
capacity, source coding efficiency, and effective information rates
across three time horizons: today (March~2026), the near term
(2027--2028), and the medium term (2029--2031). The analysis draws on
the standards, codec parameters, and infrastructure projections
developed in the main text.

\subsection{Channel Capacity of the Transport Layer}

The Shannon-Hartley theorem gives the theoretical maximum information
rate for a channel with additive white Gaussian noise:
\begin{equation}
  C = B \log_2\!\left(1 + \frac{S}{N}\right) \quad \text{[bits/s]}
  \label{eq:shannon}
\end{equation}
where $B$ is the channel bandwidth in Hz and $S/N$ is the
signal-to-noise ratio.

While the physical-layer capacity of a broadband link is governed by
(\ref{eq:shannon}), the \emph{effective} capacity available to the
application is reduced by protocol overhead, congestion control, and
transport inefficiencies. We define the effective application-layer
throughput as:
\begin{equation}
  R_{\text{eff}} = C_{\text{link}} \cdot \eta_{\text{transport}}
    \cdot (1 - \rho_{\text{overhead}})
  \label{eq:reff}
\end{equation}
where $C_{\text{link}}$ is the link-layer capacity,
$\eta_{\text{transport}} \in (0,1]$ is the transport efficiency factor
(capturing congestion control, retransmission, and multiplexing
behavior), and $\rho_{\text{overhead}} \in [0,1)$ is the fractional
protocol overhead (headers, encryption, framing).

\subsubsection{Head-of-Line Blocking and Multiplexed Capacity}

Consider $n$ independent streams multiplexed over a single connection.
Under TCP (HTTP/2), a single lost packet blocks all $n$ streams. The
expected throughput per stream under loss rate $p$ is bounded by:
\begin{equation}
  R_{\text{TCP}}^{(i)} \leq \frac{R_{\text{eff}}}{n}
    \cdot (1 - p)^{n-1}
  \label{eq:hol-tcp}
\end{equation}
because a loss in \emph{any} of the $n$ streams stalls all others. Under
QUIC (RFC~9000), streams are independent at the transport layer, so:
\begin{equation}
  R_{\text{QUIC}}^{(i)} \leq \frac{R_{\text{eff}}}{n}
    \cdot (1 - p)
  \label{eq:hol-quic}
\end{equation}
The ratio of effective per-stream throughput is:
\begin{equation}
  \frac{R_{\text{QUIC}}^{(i)}}{R_{\text{TCP}}^{(i)}}
    = \frac{1}{(1-p)^{n-2}}
  \label{eq:hol-ratio}
\end{equation}
\excursus{Excursus: Why $n = 8$?}{In a typical multi-participant video
conference, each peer contributes multiple concurrent streams: a camera video
track, an audio track, and often a screen-share or data channel. A four-person
call with two media tracks each produces $n = 8$ multiplexed streams over a
single connection --- the reference value used throughout this section. Larger
meetings (e.g.\ 10 participants with simulcast layers) push $n$ higher,
amplifying the HOL blocking effect further.}

For $n = 8$ streams (a typical multi-track video conference) at
$p = 0.01$ (1\% loss), this gives a QUIC advantage factor of
$1/(0.99)^6 \approx 1.062$, a modest 6.2\% improvement. At $p = 0.05$
(5\% loss, common on mobile or congested links), the factor is
$1/(0.95)^6 \approx 1.34$, a 34\% throughput advantage per stream.

\begin{figure}[ht]
\centering
\begin{tikzpicture}[>=Stealth]

% --- Shaded region (drawn FIRST so curves appear on top) ---
\fill[orange!8] (2, 0) rectangle (10, 5.3);
\node[font=\tiny\itshape, text=black!40, anchor=north east] at (9.8, 5.2)
  {mobile / congested};
\node[font=\tiny\itshape, text=black!40, anchor=north west] at (0.2, 5.2)
  {fixed broadband};

% --- Grid lines ---
\foreach \y in {1, 2, 3, 4, 5} {
  \draw[gray!20] (0, \y) -- (10, \y);
}

% --- Axes ---
\draw[->, thick] (0, 0) -- (10.5, 0);
\draw[->, thick] (0, 0) -- (0, 5.5);

% --- Axis titles (positioned to clear tick labels) ---
\node[font=\small] at (5, -0.85) {Packet loss rate $p$};
\node[font=\small, rotate=90] at (-0.9, 2.75) {QUIC advantage factor};

% --- X-axis labels ---
\foreach \x/\lab in {0/0, 2/0.02, 4/0.04, 6/0.06, 8/0.08, 10/0.10} {
  \draw (\x, -0.1) -- (\x, 0.1);
  \node[font=\tiny, anchor=north] at (\x, -0.15) {\lab};
}

% --- Y-axis labels ---
\foreach \y/\lab in {0/1.0, 1/1.1, 2/1.2, 3/1.3, 4/1.4, 5/1.5} {
  \draw (-0.1, \y) -- (0.1, \y);
  \node[font=\tiny, anchor=east] at (-0.15, \y) {\lab};
}

% --- n=4 curve: 1/(1-p)^2 ---
\draw[very thick, blue!70]
  (0, 0) .. controls (2.5, 0.5) and (5, 1.1) ..
  (7.5, 2.0) .. controls (8.5, 2.4) and (9.5, 2.9) ..
  (10, 3.1);
\node[font=\scriptsize, blue!70, anchor=south west] at (9.0, 3.0)
  {$n = 4$};

% --- n=8 curve: 1/(1-p)^6 ---
\draw[very thick, orange!80!black]
  (0, 0) .. controls (1.5, 0.5) and (3, 1.7) ..
  (4, 2.5) .. controls (5, 3.5) and (5.5, 4.0) ..
  (6, 4.7);
\node[font=\scriptsize, orange!80!black, anchor=south] at (5.5, 4.8)
  {$n = 8$};

% --- Reference points ---
\fill[orange!80!black] (1.0, 0.62) circle (2.5pt);
\node[font=\tiny, orange!80!black, anchor=south west] at (1.1, 0.62)
  {6.2\%};

\fill[orange!80!black] (5.0, 3.7) circle (2.5pt);
\node[font=\tiny, orange!80!black, anchor=west] at (5.1, 3.7)
  {34\%};

\fill[blue!70] (1.0, 0.20) circle (2.5pt);

% --- Annotation ---
\node[font=\scriptsize\itshape, text=black!60, align=center,
  anchor=north] at (5.0, -1.0)
  {QUIC recovers per-stream throughput closer to $C/n$ (Shannon fair share)\\
   by eliminating head-of-line blocking (eq.~\ref{eq:hol-ratio}).};

\end{tikzpicture}
\caption{QUIC advantage factor $1/(1-p)^{n-2}$ over TCP for multiplexed
  streams. At 5\% loss with $n=8$ streams, QUIC delivers 34\% more
  per-stream throughput by maintaining independence at the transport layer.}
\label{fig:hol-blocking}
\end{figure}

\subsubsection{Connection Establishment Latency}

Let $\tau$ denote the one-way latency (half the round-trip time). The
information-carrying capacity during connection establishment is zero;
the handshake represents pure overhead. The time to first useful byte
is:
\begin{align}
  T_{\text{TCP+TLS\,1.3}} &= 2\tau_{\text{TCP}} + \tau_{\text{TLS}}
    = 3\tau \quad \text{(1-RTT TLS~1.3)} \label{eq:tcp-setup} \\
  T_{\text{QUIC}} &= \tau \quad \text{(1-RTT)}
    \label{eq:quic-setup} \\
  T_{\text{QUIC,\,0-RTT}} &= 0 \quad \text{(resumption)}
    \label{eq:quic-0rtt}
\end{align}
The QUIC 0-RTT case (\ref{eq:quic-0rtt}) eliminates the handshake
entirely for resumed connections. The information lost during
handshake dead time, relative to steady-state capacity, is:
\begin{equation}
  I_{\text{lost}} = C_{\text{link}} \cdot T_{\text{setup}}
  \label{eq:ilost}
\end{equation}
For a 120~Mbps link with $\tau = 25$~ms (typical fixed broadband in
2026), the TCP+TLS~1.3 handshake wastes $120 \times 10^6 \times 0.075
= 9 \times 10^6$ bits (1.125~MB), while QUIC~1-RTT wastes only
$3 \times 10^6$ bits (375~KB).

\subsection{Source Coding Efficiency}

The rate-distortion function $R(D)$ defines the minimum bit rate
required to represent a source at distortion level $D$. For a
Gaussian source with variance $\sigma^2$ under mean squared error
distortion:
\begin{equation}
  R(D) = \frac{1}{2} \log_2 \frac{\sigma^2}{D}
    \quad \text{[bits/sample]}
  \label{eq:rate-distortion}
\end{equation}

\excursus{Excursus: Why Gaussian source?}{The Gaussian memoryless model is used
because it yields the highest (worst-case) $R(D)$ among all sources of equal
variance. Any codec that approaches the Gaussian bound on real video is
therefore doing at least as well as the bound predicts --- the true $R(D)$ for
spatially and temporally correlated video lies below the Gaussian curve, making
it a conservative benchmark.}

No practical codec achieves $R(D)$;\footnote{The $R(D)$ curve in
Figure~\ref{fig:rate-distortion} is the Gaussian memoryless source bound
(equation~\ref{eq:rate-distortion}), which is an upper bound on the true
$R(D)$ for non-Gaussian sources such as video.  Additionally, $R(D)$ assumes
zero excess-distortion probability, whereas practical codecs permit a
frame-level excess-distortion probability on the order of $10^{-6}$ to
$10^{-8}$.  Both factors allow codec operating points measured on real content
to approach or marginally cross the Gaussian bound.} real codecs operate at
some rate $R_{\text{codec}} > R(D)$. We define the coding efficiency gap as:
\begin{equation}
  \Delta = \frac{R_{\text{codec}} - R(D)}{R(D)}
  \label{eq:coding-gap}
\end{equation}

The codecs in the WebRTC stack can be ordered by their empirical
proximity to $R(D)$ for typical video content at conversational
resolutions (720p, 30~fps). Using published Bj\o ntegaard Delta (BD)
rate measurements:
\begin{equation}
  R_{\text{H.264}} > R_{\text{VP9}} > R_{\text{AV1}}
    > R(D)
  \label{eq:codec-ordering}
\end{equation}

Specifically, the empirical rate savings at equivalent PSNR are
approximately:
\begin{align}
  R_{\text{VP9}}  &\approx 0.65 \cdot R_{\text{H.264}}
    \quad (\sim\!35\% \text{ savings}) \label{eq:vp9-savings} \\
  R_{\text{AV1}}  &\approx 0.70 \cdot R_{\text{VP9}}
    \approx 0.46 \cdot R_{\text{H.264}}
    \quad (\sim\!30\% \text{ over VP9})
  \label{eq:av1-savings}
\end{align}

\begin{figure}[ht]
\centering
\begin{tikzpicture}[>=Stealth]

% --- Axes ---
\draw[->, thick] (0, 0) -- (10.5, 0);
\draw[->, thick] (0, 0) -- (0, 6.5);

% --- Axis titles ---
\node[font=\small] at (5, -0.85) {Distortion $D$};
\node[font=\small, rotate=90] at (-0.9, 3.25) {Rate [Mbps]};

% --- Y-axis rate labels ---
\foreach \y/\lab in {1.0/0.5, 2.0/1.0, 3.0/1.5, 4.0/2.0, 5.0/2.5, 6.0/3.0} {
  \draw (-0.1, \y) -- (0.1, \y);
  \node[font=\tiny, anchor=east] at (-0.15, \y) {\lab};
}

% --- R(D) Shannon bound curve ---
\draw[very thick, green!60!black]
  (1.0, 5.8) .. controls (2.0, 3.5) and (3.5, 2.0) ..
  (5.0, 1.3) .. controls (6.5, 0.8) and (8.0, 0.55) ..
  (9.5, 0.4);
\node[font=\scriptsize, green!60!black, anchor=north west] at (8.5, 0.55)
  {$R(D)$ bound};

% --- Reference distortion line (720p/30fps operating point) ---
\draw[thin, gray, dashed] (4.0, 0) -- (4.0, 6.2);
\node[font=\tiny, gray, anchor=south] at (4.0, 6.2) {720p/30fps};

% --- R(D) value at reference distortion ---
\fill[green!60!black] (4.0, 1.55) circle (2pt);

% --- Codec operating points at reference distortion ---
% H.264: ~2.5 Mbps
\draw[thick, dashed, red!70!black] (0.3, 4.8) -- (9.5, 4.8);
\fill[red!70!black] (4.0, 4.8) circle (3pt);
\node[font=\scriptsize, red!70!black, anchor=west] at (9.6, 4.8) {H.264};

% VP9: ~1.6 Mbps (0.65 * H.264)
\draw[thick, dashed, blue!70] (0.3, 3.3) -- (9.5, 3.3);
\fill[blue!70] (4.0, 3.3) circle (3pt);
\node[font=\scriptsize, blue!70, anchor=west] at (9.6, 3.3) {VP9};

% AV1: ~1.15 Mbps (0.46 * H.264)
\draw[thick, dashed, orange!80!black] (0.3, 2.3) -- (9.5, 2.3);
\fill[orange!80!black] (4.0, 2.3) circle (3pt);
\node[font=\scriptsize, orange!80!black, anchor=west] at (9.6, 2.3) {AV1};

% --- Gap annotations (braces) ---
% H.264 gap
\draw[decorate, decoration={brace, amplitude=4pt, mirror},
  thick, red!50!black]
  (3.6, 1.55) -- (3.6, 4.8)
  node[midway, anchor=east, font=\scriptsize, xshift=-5pt]
  {$\Delta_{\text{H.264}}$};

% AV1 gap
\draw[decorate, decoration={brace, amplitude=4pt},
  thick, orange!60!black]
  (4.4, 1.55) -- (4.4, 2.3)
  node[midway, anchor=west, font=\scriptsize, xshift=5pt]
  {$\Delta_{\text{AV1}}$};

% --- Annotation ---
\node[font=\scriptsize\itshape, text=black!60, align=center,
  anchor=north] at (5.0, -0.5)
  {The Shannon rate-distortion bound $R(D)$ is the theoretical minimum.\\
   AV1 narrows the gap $\Delta$ to ${\sim}0.4$ by 2030 (eq.~\ref{eq:coding-gap}).};

\end{tikzpicture}
\caption{Codec operating rates vs.\ the Shannon rate-distortion bound $R(D)$
  at 720p/30fps. Each codec operates above the theoretical minimum; the gap
  $\Delta$ (eq.~\ref{eq:coding-gap}) shrinks from H.264 through VP9 to AV1,
  approaching but never reaching the bound.}
\label{fig:rate-distortion}
\end{figure}

\subsubsection{Codec Entropy and WebCodecs}

The WebCodecs API (W3C Working Draft) exposes per-frame codec access,
allowing measurement and control of the instantaneous encoding rate.
For a video frame $f_t$ at time $t$, the encoder output has empirical
entropy:
\begin{equation}
  \hat{H}(f_t) = -\sum_{s \in \mathcal{S}}
    p(s \mid f_t) \log_2 p(s \mid f_t)
  \label{eq:frame-entropy}
\end{equation}
where $\mathcal{S}$ is the set of symbols emitted by the entropy coder
(CABAC for H.264/AV1, arithmetic coding for VP9). WebCodecs enables
applications to observe $|E(f_t)|$ (the encoded frame size in bits)
directly, which serves as a practical upper bound on
$\hat{H}(f_t)$:
\begin{equation}
  \hat{H}(f_t) \leq |E(f_t)| \leq R_{\text{codec}} / \text{fps}
  \label{eq:entropy-bound}
\end{equation}

\subsection{Effective Information Rate by Time Horizon}

We now combine transport and source coding measures to estimate the
effective information rate for a representative real-time video
session: a single 720p/30fps video call with one audio channel.

\excursus{Excursus: Why 720p?}{Throughout this appendix, 720p/30fps is used as
a fixed reference for cross-era comparison, not as a prediction that 720p
remains the dominant conferencing resolution. The bandwidth surplus quantified
below is precisely what enables migration to 1080p (feasible by ${\sim}$2028
with AV1 hardware encode at today's 720p bit rates) and 4K (feasible by
${\sim}$2030). The 720p baseline makes the efficiency gains legible: the same
task gets cheaper, freeing capacity for higher quality or more streams.}

Define the \emph{net information rate} as:
\begin{equation}
  I_{\text{net}} = R_{\text{codec}} \cdot \eta_{\text{transport}}
    \cdot (1 - \rho_{\text{overhead}})
    \cdot (1 - p_{\text{loss}})
  \label{eq:inet}
\end{equation}

\begin{table}[ht]
\centering
\small
\renewcommand{\arraystretch}{1.4}
\begin{tabularx}{\textwidth}{lXXX}
\toprule
\textbf{Parameter} & \textbf{March 2026} & \textbf{2027--2028} &
  \textbf{2029--2031} \\
\midrule
Median link capacity &
  120~Mbps &
  $\sim$150~Mbps &
  200--300~Mbps \\
Dominant codec (WebRTC) &
  VP9 / H.264 &
  VP9 / AV1 (HW) &
  AV1 (HW default) \\
Codec rate (720p/30fps) &
  1.5--2.5~Mbps &
  1.0--1.8~Mbps &
  0.7--1.2~Mbps \\
$\eta_{\text{transport}}$ &
  0.92 (QUIC) &
  0.93 &
  0.94 \\
$\rho_{\text{overhead}}$ (SRTP+QUIC) &
  0.06 &
  0.05 &
  0.04 \\
Typical loss $p$ &
  0.01--0.03 &
  0.01--0.02 &
  0.005--0.01 \\
$I_{\text{net}}$ (720p video) &
  1.3--2.2~Mbps &
  0.9--1.6~Mbps &
  0.6--1.1~Mbps \\
$I_{\text{net}}$ as \% of $C_{\text{link}}$ &
  1.1--1.8\% &
  0.6--1.1\% &
  0.2--0.6\% \\
\bottomrule
\end{tabularx}
\caption{Effective information rate for a 720p/30fps WebRTC video stream.
  The declining ratio $I_{\text{net}} / C_{\text{link}}$ reflects both
  improved codec efficiency and growing link capacity, freeing bandwidth
  for higher resolutions, additional streams, or concurrent data
  channels.}
\label{tab:inet}
\end{table}

The declining ratio $I_{\text{net}} / C_{\text{link}}$ in
Table~\ref{tab:inet} is the central quantitative observation: by 2030, a
720p video call consumes less than 0.5\% of available link capacity. The
surplus capacity enables qualitative changes --- 4K video conferencing,
multi-stream layouts, or concurrent data channels via WebTransport ---
that are quantitatively infeasible in 2026.

\begin{figure}[ht]
\centering
\begin{tikzpicture}[>=Stealth]

% --- Axes ---
\draw[->, thick] (0, 0) -- (11, 0);
\draw[->, thick] (0, 0) -- (0, 5.5);

% --- Axis title ---
\node[font=\small, rotate=90] at (-0.9, 2.75) {Mbps};

% --- Y-axis labels ---
\foreach \y/\lab in {0/0, 1/50, 2/100, 3/150, 4/200, 5/250} {
  \draw (-0.1, \y) -- (0.1, \y);
  \node[font=\tiny, anchor=east] at (-0.15, \y) {\lab};
}

% --- Group: March 2026 ---
% C_link = 120 Mbps -> y = 2.4
\fill[blue!12] (0.8, 0) rectangle (2.4, 2.4);
\draw[blue!50, thick] (0.8, 0) rectangle (2.4, 2.4);
% I_net = 1.75 Mbps (midpoint) -> y = 0.035
\fill[orange!60] (0.8, 0) rectangle (2.4, 0.07);
\draw[orange!80!black, thick] (0.8, 0) rectangle (2.4, 0.07);
\node[font=\footnotesize, anchor=south] at (1.6, 2.5) {120};
\node[font=\tiny, orange!80!black] at (1.6, -0.4) {1.5\%};
\node[font=\small, anchor=north] at (1.6, -0.7) {2026};

% --- Group: 2027--2028 ---
% C_link = 150 Mbps -> y = 3.0
\fill[blue!12] (4.0, 0) rectangle (5.6, 3.0);
\draw[blue!50, thick] (4.0, 0) rectangle (5.6, 3.0);
% I_net = 1.25 Mbps (midpoint) -> y = 0.025
\fill[orange!60] (4.0, 0) rectangle (5.6, 0.05);
\draw[orange!80!black, thick] (4.0, 0) rectangle (5.6, 0.05);
\node[font=\footnotesize, anchor=south] at (4.8, 3.1) {150};
\node[font=\tiny, orange!80!black] at (4.8, -0.4) {0.8\%};
\node[font=\small, anchor=north] at (4.8, -0.7) {2027--28};

% --- Group: 2029--2031 ---
% C_link = 250 Mbps -> y = 5.0
\fill[blue!12] (7.2, 0) rectangle (8.8, 5.0);
\draw[blue!50, thick] (7.2, 0) rectangle (8.8, 5.0);
% I_net = 0.85 Mbps (midpoint) -> y = 0.017
\fill[orange!60] (7.2, 0) rectangle (8.8, 0.034);
\draw[orange!80!black, thick] (7.2, 0) rectangle (8.8, 0.034);
\node[font=\footnotesize, anchor=south] at (8.0, 5.1) {250};
\node[font=\tiny, orange!80!black] at (8.0, -0.4) {0.3\%};
\node[font=\small, anchor=north] at (8.0, -0.7) {2029--31};

% --- Legend ---
\node[anchor=west, font=\footnotesize] at (0.5, 5.8) {%
  \tikz[baseline=-0.5ex]{\fill[blue!12] (0,0) rectangle (0.3,0.3);
    \draw[blue!50, thick] (0,0) rectangle (0.3,0.3);}
  ~$C_{\text{link}}$ (channel capacity) \qquad
  \tikz[baseline=-0.5ex]{\fill[orange!60] (0,0) rectangle (0.3,0.3);
    \draw[orange!80!black, thick] (0,0) rectangle (0.3,0.3);}
  ~$I_{\text{net}}$ (720p stream)};

% --- Annotation ---
\node[font=\scriptsize\itshape, text=black!60, align=center,
  anchor=north] at (5.5, -1.2)
  {A single 720p call uses a vanishing fraction of Shannon-derived\\
   channel capacity, freeing bandwidth for 4K, multi-stream, and data channels.};

\end{tikzpicture}
\caption{Channel capacity $C_{\text{link}}$ vs.\ effective information rate
  $I_{\text{net}}$ for a single 720p/30fps stream across time horizons. The
  orange bar (stream rate) is barely visible against the blue bar (capacity),
  illustrating the growing surplus as codecs improve and bandwidth increases.}
\label{fig:inet-capacity}
\end{figure}

\subsection{Security Entropy and Key Space}

The DTLS-SRTP security architecture (RFCs~8827, 3711) protects all
WebRTC media. Information-theoretic security requires that the key
entropy exceed the information content of the protected material.

For a DTLS~1.3 handshake with ECDHE on P-256, the key agreement
provides:
\begin{equation}
  H_{\text{key}} = 128 \text{~bits of security}
  \label{eq:key-entropy}
\end{equation}

The SRTP master key is 128 bits (AES-128-GCM), giving an exhaustive
search cost of $2^{128}$ operations. The total information content of a
one-hour 720p video call at 2~Mbps is:
\begin{equation}
  I_{\text{call}} = 2 \times 10^6 \times 3600
    = 7.2 \times 10^9 \text{~bits}
    \approx 2^{32.7} \text{~bits}
  \label{eq:call-content}
\end{equation}

The ratio of key entropy to protected content:
\begin{equation}
  \frac{H_{\text{key}}}{I_{\text{call}}}
    = \frac{128}{32.7} \approx 3.91
  \label{eq:security-margin}
\end{equation}

This ratio remains comfortably above~1 even for much longer sessions.
For a continuous 24-hour 4K AV1 stream at 15~Mbps:
\begin{equation}
  I_{\text{stream}} = 15 \times 10^6 \times 86400
    = 1.296 \times 10^{12} \text{~bits}
    \approx 2^{40.2} \text{~bits}
\end{equation}
yielding $H_{\text{key}} / I_{\text{stream}} \approx 3.18$, still
well above the threshold.

\subsubsection{End-to-End Encryption via Encoded Transforms}

The RTCRtpScriptTransform API (Encoded Transforms) enables
application-layer E2EE within the WebRTC pipeline. From an
information-theoretic perspective, this introduces a second encryption
layer with its own key, creating a \emph{nested channel}:
\begin{equation}
  X \xrightarrow{\;K_{\text{E2EE}}\;}
  E_1(X) \xrightarrow{\;K_{\text{SRTP}}\;}
  E_2(E_1(X))
  \label{eq:nested}
\end{equation}

The mutual information between the ciphertext and plaintext, given
neither key, satisfies:
\begin{equation}
  I\!\left(X;\, E_2(E_1(X))\right) \leq
    \min\!\big(H(K_{\text{E2EE}})^{-1},\;
    H(K_{\text{SRTP}})^{-1}\big)
    \cdot I_{\text{call}}
  \label{eq:mutual-e2ee}
\end{equation}

In practice, with $H(K_{\text{E2EE}}) = H(K_{\text{SRTP}}) = 128$
bits, this quantity is negligible ($\ll 1$ bit for any practical
session length), confirming that the nested encryption provides
information-theoretic confidentiality well beyond computational
security guarantees.

Today (March~2026), Encoded Transforms remain Chrome-only. By
2027--2028 (Section~8.3), cross-browser standardization will make
this nested channel universally available. By 2029--2031, SFrame
for group calls will extend this to $n$-party sessions, where the
mutual information between any non-member and the group plaintext
remains negligible:
\begin{equation}
  I\!\left(X;\, E(X) \;\middle|\; K_i \notin
    \{K_1, \ldots, K_n\}\right) \approx 0
  \label{eq:sframe}
\end{equation}

\subsection{WebTransport: Unreliable Datagrams and Channel Capacity}

WebTransport (RFC~9297) introduces unreliable datagrams to browser-server
communication. For latency-sensitive applications (game state, sensor
telemetry), the relevant measure is not throughput but
\emph{freshness-weighted information rate}. Define the value of a
datagram as exponentially decaying with delivery latency $\delta$:
\begin{equation}
  V(\delta) = e^{-\lambda \delta}
  \label{eq:freshness}
\end{equation}
where $\lambda > 0$ is an application-specific decay rate. The
effective information rate under unreliable delivery is:
\begin{equation}
  I_{\text{WT}} = R_{\text{send}} \cdot (1 - p_{\text{loss}})
    \cdot \mathbb{E}\!\left[e^{-\lambda \delta}\right]
  \label{eq:iwt}
\end{equation}

Under reliable delivery (TCP/WebSocket), a retransmission delays
subsequent data, so:
\begin{equation}
  I_{\text{WS}} = R_{\text{send}}
    \cdot \mathbb{E}\!\left[e^{-\lambda(\delta + p \cdot 2\tau)}\right]
  \label{eq:iws}
\end{equation}

The ratio $I_{\text{WT}} / I_{\text{WS}}$ exceeds 1 when:
\begin{equation}
  (1 - p) \cdot e^{\lambda \cdot p \cdot 2\tau} > 1
  \quad \Longleftrightarrow \quad
  \lambda > \frac{-\ln(1 - p)}{2 p \tau}
  \label{eq:wt-advantage}
\end{equation}

For $p = 0.02$ and $\tau = 25$~ms, this gives $\lambda > 20.2$~s$^{-1}$
(half-life $< 34$~ms).

\excursus{Excursus: Why $p = 0.02$ and $\tau = 25$\,ms?}{These values represent
median fixed-broadband conditions in 2026 (Table~\ref{tab:inet}). On mobile or
congested links ($p \approx 0.05$, $\tau \approx 50$\,ms), the threshold
$\lambda^*$ drops to ${\sim}10$\,s$^{-1}$, broadening WebTransport's advantage
to applications with data half-lives up to ${\sim}70$\,ms. The fixed-broadband
reference is thus conservative: mobile scenarios favor unreliable delivery even
more strongly.}

Any application where data older than
$\sim$30~ms is useless benefits from WebTransport's unreliable
mode over WebSocket --- precisely the regime of interactive gaming, live
cursor tracking, and real-time sensor feeds.

As of March~2026, Safari does not support WebTransport
(Section~5), limiting coverage to $\sim$82\%. By 2027--2028
(Section~8.2), universal support will allow applications to rely on
(\ref{eq:iwt}) without fallback paths.

\begin{figure}[ht]
\centering
\begin{tikzpicture}[>=Stealth]

% --- Axes ---
\draw[->, thick] (0, 0) -- (10.5, 0);
\draw[->, thick] (0, 0) -- (0, 5.5);

% --- Axis titles (positioned to clear tick labels) ---
\node[font=\small] at (5, -0.75)
  {Freshness decay rate $\lambda$ [s$^{-1}$]};
\node[font=\small, rotate=90] at (-1.1, 2.75)
  {$I_{\text{WT}} / I_{\text{WS}}$};

% --- X-axis labels ---
\foreach \x/\lab in {0/0, 1.67/10, 3.33/20, 5.0/30, 6.67/40, 8.33/50, 10/60} {
  \draw (\x, -0.1) -- (\x, 0.1);
  \node[font=\tiny, anchor=north] at (\x, -0.15) {\lab};
}

% --- Y-axis labels (0.96 to 1.06, each 0.02 = 1 unit) ---
\foreach \y/\lab in {0/0.96, 1/0.98, 2/1.00, 3/1.02, 4/1.04, 5/1.06} {
  \draw (-0.1, \y) -- (0.1, \y);
  \node[font=\tiny, anchor=east] at (-0.15, \y) {\lab};
}

% --- Ratio = 1.0 line (y=2 maps to ratio 1.0) ---
\draw[thin, gray, dashed] (0, 2) -- (10, 2);

% --- Shade WebTransport advantage region ---
% lambda* = 20.2 -> x = 3.37
\fill[green!10] (3.37, 2) rectangle (10, 5.3);

% --- I_WT/I_WS curve: 0.98*exp(lambda*0.001) ---
% Y range: 0.96-1.06 (each 0.02 = 1 unit). lambda=0: y=1, lambda*=20.2: y=2, lambda=60: y=4.05
\draw[very thick, orange!80!black]
  (0, 1.0) .. controls (1.5, 1.3) and (2.5, 1.7) ..
  (3.37, 2.0) .. controls (4.5, 2.4) and (6, 2.8) ..
  (8, 3.4) .. controls (9, 3.7) and (9.5, 3.85) ..
  (10, 4.05);

% --- Threshold line ---
\draw[thick, dashed, red!60!black] (3.37, 0) -- (3.37, 5.2);
\node[font=\scriptsize, red!60!black, anchor=south, rotate=90] at (3.17, 3.5)
  {$\lambda^* = 20.2$ s$^{-1}$};

% --- Application regime labels ---
\node[font=\scriptsize\itshape, text=blue!60, anchor=north] at (1.5, 5.3)
  {video conf};
\node[font=\scriptsize\itshape, text=green!50!black, anchor=north] at (7.0, 5.3)
  {gaming / cursor / sensors};

% --- Mark advantage/disadvantage ---
\node[font=\tiny, text=red!50!black, anchor=north east] at (3.2, 1.8)
  {WebSocket wins};
\node[font=\tiny, text=green!50!black, anchor=north west] at (3.5, 4.5)
  {WebTransport wins};

% --- Annotation ---
\node[font=\scriptsize\itshape, text=black!60, align=center,
  anchor=north] at (5.0, -1.0)
  {For $p = 0.02$, $\tau = 25$~ms. Applications where data older than\\
   ${\sim}$30~ms is stale gain from unreliable delivery (eq.~\ref{eq:wt-advantage}).};

\end{tikzpicture}
\caption{Freshness-weighted information rate ratio $I_{\text{WT}} / I_{\text{WS}}$
  as a function of decay rate $\lambda$. Above the threshold $\lambda^*$,
  WebTransport's unreliable datagrams deliver more useful information per
  unit time than WebSocket's reliable delivery, because retransmission
  wastes capacity on stale data.}
\label{fig:freshness}
\end{figure}

\subsection{Projections Summary and Practical Implications}

\begin{table}[ht]
\centering
\small
\renewcommand{\arraystretch}{1.4}
\begin{tabularx}{\textwidth}{lXXX}
\toprule
\textbf{Measure} & \textbf{March 2026} & \textbf{2027--2028} &
  \textbf{2029--2031} \\
\midrule
Channel capacity $C_{\text{link}}$ &
  120~Mbps &
  150~Mbps &
  250~Mbps \\
HOL blocking factor $(1\!-\!p)^{n-2}$ at $n\!=\!8$ &
  0.94 ($p\!=\!0.01$) &
  0.94 &
  0.97 ($p\!=\!0.005$) \\
Source coding gap $\Delta$ &
  VP9: $\sim$0.8 &
  AV1 HW: $\sim$0.5 &
  AV1 opt: $\sim$0.4 \\
$I_{\text{net}}$ per 720p stream &
  1.3--2.2~Mbps &
  0.9--1.6~Mbps &
  0.6--1.1~Mbps \\
Max 720p streams in $C_{\text{link}}$ &
  $\sim$55--90 &
  $\sim$95--165 &
  $\sim$230--415 \\
E2EE availability &
  Chrome only &
  Cross-browser &
  SFrame $n$-party \\
WebTransport coverage &
  82\% &
  $\sim$100\% &
  100\% \\
Freshness threshold $\lambda^*$ &
  20.2~s$^{-1}$ &
  $\sim$18~s$^{-1}$ &
  $\sim$15~s$^{-1}$ \\
\bottomrule
\end{tabularx}
\caption{Information-theoretic measures across time horizons. $\Delta$
  is the coding efficiency gap (\ref{eq:coding-gap}), $I_{\text{net}}$
  is the net information rate (\ref{eq:inet}), and $\lambda^*$ is the
  freshness decay threshold above which WebTransport dominates WebSocket
  (\ref{eq:wt-advantage}).}
\label{tab:projections}
\end{table}

The trajectory is clear: the combination of improving source coding
(AV1 hardware), expanding channel capacity (broadband growth), and
maturing transport protocols (QUIC universal, WebTransport universal)
drives the \emph{information cost per unit of communicated experience}
steadily downward. The remaining question is not whether the
Shannon capacity suffices --- it already does, by orders of magnitude
--- but how efficiently the protocol stack converts that capacity into
delivered, fresh, secure information at the application layer. By
2030, the gap between theoretical capacity and realized throughput will
be the narrowest it has ever been for browser-based real-time
communication.

\subsubsection{What the Theory Means for Applications}

The equations and figures in this appendix are not purely academic. Each
theoretical result maps to a concrete design decision that application
developers face today and over the next five years.

\textbf{Codec selection determines bandwidth cost, not quality.}
The rate-distortion analysis (Section~A.2, Figure~\ref{fig:rate-distortion})
shows that AV1 operates at roughly 46\% of H.264's bit rate for equivalent
perceptual quality at 720p. In practice, this means a WebRTC application
switching from H.264 to AV1 can either halve its bandwidth consumption at
the same quality or substantially increase resolution (e.g., 720p to 1080p)
within the same bandwidth envelope. By 2030, with AV1 hardware encoding
ubiquitous and the coding gap $\Delta$ approaching 0.4, further codec
improvements yield diminishing returns --- the stack will be operating close
enough to the Shannon rate-distortion bound that the next generation of
efficiency gains will come from transport and architecture, not codecs alone.

\textbf{QUIC's advantage is situational but decisive on bad networks.}
The head-of-line blocking analysis (Section~A.1.1,
Figure~\ref{fig:hol-blocking}) shows that QUIC's per-stream throughput
advantage over TCP is modest on good networks (6\% at 1\% loss) but
substantial on degraded ones (34\% at 5\% loss). The practical implication
is that applications targeting mobile users, emerging markets, or congested
enterprise networks --- precisely the environments where real-time
communication is hardest --- benefit disproportionately from QUIC-based
transports. For a video conference with 8 participants on a cellular
connection experiencing 3--5\% packet loss, the difference between QUIC and
TCP is the difference between usable and degraded video.

\textbf{Connection latency is a solved problem for returning users.}
The handshake analysis (Section~A.1.2) shows that QUIC 0-RTT eliminates
connection establishment overhead entirely for resumed connections. For a
WebTransport application that a user opens daily --- a collaboration tool,
a live dashboard, a recurring meeting client --- the time to first useful
byte drops to zero. The 1.125~MB of wasted capacity during a TCP+TLS
handshake on a 120~Mbps link may seem small, but for latency-sensitive
applications where the first 100~ms of perceived responsiveness determines
user experience, the difference is measurable.

\textbf{Bandwidth surplus enables architectural ambition.}
The effective information rate analysis (Section~A.3,
Figure~\ref{fig:inet-capacity}) reveals the most consequential practical
finding: a single 720p video stream consumes a vanishing fraction of
available link capacity (under 0.5\% by 2030). This surplus is not merely
``headroom'' --- it is the resource that makes new application architectures
feasible. A 250~Mbps link can theoretically carry 230--415 simultaneous 720p
streams. In practice, this means applications can afford to run multiple
concurrent media flows (video, screen share, spatial audio), data channels
(collaborative editing, file transfer), and telemetry streams without
contending for bandwidth. The constraint shifts from ``can the network carry
this?''\ to ``can the application manage this complexity?''

\textbf{End-to-end encryption adds negligible overhead.}
The security analysis (Section~A.4) confirms that the nested DTLS-SRTP plus
E2EE architecture imposes no meaningful information-theoretic cost. The key
entropy ($2^{128}$) exceeds the information content of any practical session
by a factor of at least 3. For application developers, this means E2EE can
be treated as a default rather than a premium feature --- there is no
bandwidth, latency, or capacity penalty for encrypting at the application
layer on top of the transport-layer encryption that WebRTC already mandates.
The only constraint is API availability, which moves from Chrome-only (2026)
to universal (2028).

\textbf{WebTransport unlocks real-time data that WebSocket cannot serve.}
The freshness analysis (Section~A.5, Figure~\ref{fig:freshness}) quantifies
when unreliable delivery outperforms reliable delivery: any data with a
useful half-life below approximately 34~ms benefits from WebTransport's
unreliable datagrams. This covers a wide class of interactive applications
--- game state synchronization, live cursor and pointer sharing, real-time
sensor feeds, collaborative whiteboard strokes --- where retransmitting a
stale update is worse than dropping it. Before WebTransport, these
applications had to choose between WebSocket (reliable but stale under loss)
and WebRTC DataChannel (unreliable mode available but requiring the full
peer-to-peer ICE/DTLS machinery). WebTransport provides the unreliable
mode with a simple client-server model over HTTP/3, removing the
architectural overhead.

\textbf{The overall picture: the stack converges on efficiency.}
Taken together, the information-theoretic measures trace a consistent
trajectory. The protocol stack is converging toward the theoretical limits
established by Shannon: codecs approach the rate-distortion bound, transports
eliminate unnecessary blocking and handshake overhead, and encryption adds
negligible cost. By 2030, the practical distance between what the standards
deliver and what information theory permits will be the smallest it has ever
been for browser-based real-time communication. The consequence for
application developers is that the infrastructure --- standards, hardware,
and networks --- will no longer be the binding constraint. The remaining
challenge is purely architectural: how to compose these efficient, universal
primitives into applications that fully exploit the capacity and low latency
they provide.

%
% Requires in the preamble: amsmath, amssymb, tabularx, booktabs

\section{Information-Theoretic View}
\label{app:info-theory}

This appendix applies information-theoretic measures to the real-time
web communication stack described in this document. We quantify channel
capacity, source coding efficiency, and effective information rates
across three time horizons: today (March~2026), the near term
(2027--2028), and the medium term (2029--2031). The analysis draws on
the standards, codec parameters, and infrastructure projections
developed in the main text.

\subsection{Channel Capacity of the Transport Layer}

The Shannon-Hartley theorem gives the theoretical maximum information
rate for a channel with additive white Gaussian noise:
\begin{equation}
  C = B \log_2\!\left(1 + \frac{S}{N}\right) \quad \text{[bits/s]}
  \label{eq:shannon}
\end{equation}
where $B$ is the channel bandwidth in Hz and $S/N$ is the
signal-to-noise ratio.

While the physical-layer capacity of a broadband link is governed by
(\ref{eq:shannon}), the \emph{effective} capacity available to the
application is reduced by protocol overhead, congestion control, and
transport inefficiencies. We define the effective application-layer
throughput as:
\begin{equation}
  R_{\text{eff}} = C_{\text{link}} \cdot \eta_{\text{transport}}
    \cdot (1 - \rho_{\text{overhead}})
  \label{eq:reff}
\end{equation}
where $C_{\text{link}}$ is the link-layer capacity,
$\eta_{\text{transport}} \in (0,1]$ is the transport efficiency factor
(capturing congestion control, retransmission, and multiplexing
behavior), and $\rho_{\text{overhead}} \in [0,1)$ is the fractional
protocol overhead (headers, encryption, framing).

\subsubsection{Head-of-Line Blocking and Multiplexed Capacity}

Consider $n$ independent streams multiplexed over a single connection.
Under TCP (HTTP/2), a single lost packet blocks all $n$ streams. The
expected throughput per stream under loss rate $p$ is bounded by:
\begin{equation}
  R_{\text{TCP}}^{(i)} \leq \frac{R_{\text{eff}}}{n}
    \cdot (1 - p)^{n-1}
  \label{eq:hol-tcp}
\end{equation}
because a loss in \emph{any} of the $n$ streams stalls all others. Under
QUIC (RFC~9000), streams are independent at the transport layer, so:
\begin{equation}
  R_{\text{QUIC}}^{(i)} \leq \frac{R_{\text{eff}}}{n}
    \cdot (1 - p)
  \label{eq:hol-quic}
\end{equation}
The ratio of effective per-stream throughput is:
\begin{equation}
  \frac{R_{\text{QUIC}}^{(i)}}{R_{\text{TCP}}^{(i)}}
    = \frac{1}{(1-p)^{n-2}}
  \label{eq:hol-ratio}
\end{equation}
\excursus{Excursus: Why $n = 8$?}{In a typical multi-participant video
conference, each peer contributes multiple concurrent streams: a camera video
track, an audio track, and often a screen-share or data channel. A four-person
call with two media tracks each produces $n = 8$ multiplexed streams over a
single connection --- the reference value used throughout this section. Larger
meetings (e.g.\ 10 participants with simulcast layers) push $n$ higher,
amplifying the HOL blocking effect further.}

For $n = 8$ streams (a typical multi-track video conference) at
$p = 0.01$ (1\% loss), this gives a QUIC advantage factor of
$1/(0.99)^6 \approx 1.062$, a modest 6.2\% improvement. At $p = 0.05$
(5\% loss, common on mobile or congested links), the factor is
$1/(0.95)^6 \approx 1.34$, a 34\% throughput advantage per stream.

\begin{figure}[ht]
\centering
\begin{tikzpicture}[>=Stealth]

% --- Shaded region (drawn FIRST so curves appear on top) ---
\fill[orange!8] (2, 0) rectangle (10, 5.3);
\node[font=\tiny\itshape, text=black!40, anchor=north east] at (9.8, 5.2)
  {mobile / congested};
\node[font=\tiny\itshape, text=black!40, anchor=north west] at (0.2, 5.2)
  {fixed broadband};

% --- Grid lines ---
\foreach \y in {1, 2, 3, 4, 5} {
  \draw[gray!20] (0, \y) -- (10, \y);
}

% --- Axes ---
\draw[->, thick] (0, 0) -- (10.5, 0);
\draw[->, thick] (0, 0) -- (0, 5.5);

% --- Axis titles (positioned to clear tick labels) ---
\node[font=\small] at (5, -0.85) {Packet loss rate $p$};
\node[font=\small, rotate=90] at (-0.9, 2.75) {QUIC advantage factor};

% --- X-axis labels ---
\foreach \x/\lab in {0/0, 2/0.02, 4/0.04, 6/0.06, 8/0.08, 10/0.10} {
  \draw (\x, -0.1) -- (\x, 0.1);
  \node[font=\tiny, anchor=north] at (\x, -0.15) {\lab};
}

% --- Y-axis labels ---
\foreach \y/\lab in {0/1.0, 1/1.1, 2/1.2, 3/1.3, 4/1.4, 5/1.5} {
  \draw (-0.1, \y) -- (0.1, \y);
  \node[font=\tiny, anchor=east] at (-0.15, \y) {\lab};
}

% --- n=4 curve: 1/(1-p)^2 ---
\draw[very thick, blue!70]
  (0, 0) .. controls (2.5, 0.5) and (5, 1.1) ..
  (7.5, 2.0) .. controls (8.5, 2.4) and (9.5, 2.9) ..
  (10, 3.1);
\node[font=\scriptsize, blue!70, anchor=south west] at (9.0, 3.0)
  {$n = 4$};

% --- n=8 curve: 1/(1-p)^6 ---
\draw[very thick, orange!80!black]
  (0, 0) .. controls (1.5, 0.5) and (3, 1.7) ..
  (4, 2.5) .. controls (5, 3.5) and (5.5, 4.0) ..
  (6, 4.7);
\node[font=\scriptsize, orange!80!black, anchor=south] at (5.5, 4.8)
  {$n = 8$};

% --- Reference points ---
\fill[orange!80!black] (1.0, 0.62) circle (2.5pt);
\node[font=\tiny, orange!80!black, anchor=south west] at (1.1, 0.62)
  {6.2\%};

\fill[orange!80!black] (5.0, 3.7) circle (2.5pt);
\node[font=\tiny, orange!80!black, anchor=west] at (5.1, 3.7)
  {34\%};

\fill[blue!70] (1.0, 0.20) circle (2.5pt);

% --- Annotation ---
\node[font=\scriptsize\itshape, text=black!60, align=center,
  anchor=north] at (5.0, -1.0)
  {QUIC recovers per-stream throughput closer to $C/n$ (Shannon fair share)\\
   by eliminating head-of-line blocking (eq.~\ref{eq:hol-ratio}).};

\end{tikzpicture}
\caption{QUIC advantage factor $1/(1-p)^{n-2}$ over TCP for multiplexed
  streams. At 5\% loss with $n=8$ streams, QUIC delivers 34\% more
  per-stream throughput by maintaining independence at the transport layer.}
\label{fig:hol-blocking}
\end{figure}

\subsubsection{Connection Establishment Latency}

Let $\tau$ denote the one-way latency (half the round-trip time). The
information-carrying capacity during connection establishment is zero;
the handshake represents pure overhead. The time to first useful byte
is:
\begin{align}
  T_{\text{TCP+TLS\,1.3}} &= 2\tau_{\text{TCP}} + \tau_{\text{TLS}}
    = 3\tau \quad \text{(1-RTT TLS~1.3)} \label{eq:tcp-setup} \\
  T_{\text{QUIC}} &= \tau \quad \text{(1-RTT)}
    \label{eq:quic-setup} \\
  T_{\text{QUIC,\,0-RTT}} &= 0 \quad \text{(resumption)}
    \label{eq:quic-0rtt}
\end{align}
The QUIC 0-RTT case (\ref{eq:quic-0rtt}) eliminates the handshake
entirely for resumed connections. The information lost during
handshake dead time, relative to steady-state capacity, is:
\begin{equation}
  I_{\text{lost}} = C_{\text{link}} \cdot T_{\text{setup}}
  \label{eq:ilost}
\end{equation}
For a 120~Mbps link with $\tau = 25$~ms (typical fixed broadband in
2026), the TCP+TLS~1.3 handshake wastes $120 \times 10^6 \times 0.075
= 9 \times 10^6$ bits (1.125~MB), while QUIC~1-RTT wastes only
$3 \times 10^6$ bits (375~KB).

\subsection{Source Coding Efficiency}

The rate-distortion function $R(D)$ defines the minimum bit rate
required to represent a source at distortion level $D$. For a
Gaussian source with variance $\sigma^2$ under mean squared error
distortion:
\begin{equation}
  R(D) = \frac{1}{2} \log_2 \frac{\sigma^2}{D}
    \quad \text{[bits/sample]}
  \label{eq:rate-distortion}
\end{equation}

\excursus{Excursus: Why Gaussian source?}{The Gaussian memoryless model is used
because it yields the highest (worst-case) $R(D)$ among all sources of equal
variance. Any codec that approaches the Gaussian bound on real video is
therefore doing at least as well as the bound predicts --- the true $R(D)$ for
spatially and temporally correlated video lies below the Gaussian curve, making
it a conservative benchmark.}

No practical codec achieves $R(D)$;\footnote{The $R(D)$ curve in
Figure~\ref{fig:rate-distortion} is the Gaussian memoryless source bound
(equation~\ref{eq:rate-distortion}), which is an upper bound on the true
$R(D)$ for non-Gaussian sources such as video.  Additionally, $R(D)$ assumes
zero excess-distortion probability, whereas practical codecs permit a
frame-level excess-distortion probability on the order of $10^{-6}$ to
$10^{-8}$.  Both factors allow codec operating points measured on real content
to approach or marginally cross the Gaussian bound.} real codecs operate at
some rate $R_{\text{codec}} > R(D)$. We define the coding efficiency gap as:
\begin{equation}
  \Delta = \frac{R_{\text{codec}} - R(D)}{R(D)}
  \label{eq:coding-gap}
\end{equation}

The codecs in the WebRTC stack can be ordered by their empirical
proximity to $R(D)$ for typical video content at conversational
resolutions (720p, 30~fps). Using published Bj\o ntegaard Delta (BD)
rate measurements:
\begin{equation}
  R_{\text{H.264}} > R_{\text{VP9}} > R_{\text{AV1}}
    > R(D)
  \label{eq:codec-ordering}
\end{equation}

Specifically, the empirical rate savings at equivalent PSNR are
approximately:
\begin{align}
  R_{\text{VP9}}  &\approx 0.65 \cdot R_{\text{H.264}}
    \quad (\sim\!35\% \text{ savings}) \label{eq:vp9-savings} \\
  R_{\text{AV1}}  &\approx 0.70 \cdot R_{\text{VP9}}
    \approx 0.46 \cdot R_{\text{H.264}}
    \quad (\sim\!30\% \text{ over VP9})
  \label{eq:av1-savings}
\end{align}

\begin{figure}[ht]
\centering
\begin{tikzpicture}[>=Stealth]

% --- Axes ---
\draw[->, thick] (0, 0) -- (10.5, 0);
\draw[->, thick] (0, 0) -- (0, 6.5);

% --- Axis titles ---
\node[font=\small] at (5, -0.85) {Distortion $D$};
\node[font=\small, rotate=90] at (-0.9, 3.25) {Rate [Mbps]};

% --- Y-axis rate labels ---
\foreach \y/\lab in {1.0/0.5, 2.0/1.0, 3.0/1.5, 4.0/2.0, 5.0/2.5, 6.0/3.0} {
  \draw (-0.1, \y) -- (0.1, \y);
  \node[font=\tiny, anchor=east] at (-0.15, \y) {\lab};
}

% --- R(D) Shannon bound curve ---
\draw[very thick, green!60!black]
  (1.0, 5.8) .. controls (2.0, 3.5) and (3.5, 2.0) ..
  (5.0, 1.3) .. controls (6.5, 0.8) and (8.0, 0.55) ..
  (9.5, 0.4);
\node[font=\scriptsize, green!60!black, anchor=north west] at (8.5, 0.55)
  {$R(D)$ bound};

% --- Reference distortion line (720p/30fps operating point) ---
\draw[thin, gray, dashed] (4.0, 0) -- (4.0, 6.2);
\node[font=\tiny, gray, anchor=south] at (4.0, 6.2) {720p/30fps};

% --- R(D) value at reference distortion ---
\fill[green!60!black] (4.0, 1.55) circle (2pt);

% --- Codec operating points at reference distortion ---
% H.264: ~2.5 Mbps
\draw[thick, dashed, red!70!black] (0.3, 4.8) -- (9.5, 4.8);
\fill[red!70!black] (4.0, 4.8) circle (3pt);
\node[font=\scriptsize, red!70!black, anchor=west] at (9.6, 4.8) {H.264};

% VP9: ~1.6 Mbps (0.65 * H.264)
\draw[thick, dashed, blue!70] (0.3, 3.3) -- (9.5, 3.3);
\fill[blue!70] (4.0, 3.3) circle (3pt);
\node[font=\scriptsize, blue!70, anchor=west] at (9.6, 3.3) {VP9};

% AV1: ~1.15 Mbps (0.46 * H.264)
\draw[thick, dashed, orange!80!black] (0.3, 2.3) -- (9.5, 2.3);
\fill[orange!80!black] (4.0, 2.3) circle (3pt);
\node[font=\scriptsize, orange!80!black, anchor=west] at (9.6, 2.3) {AV1};

% --- Gap annotations (braces) ---
% H.264 gap
\draw[decorate, decoration={brace, amplitude=4pt, mirror},
  thick, red!50!black]
  (3.6, 1.55) -- (3.6, 4.8)
  node[midway, anchor=east, font=\scriptsize, xshift=-5pt]
  {$\Delta_{\text{H.264}}$};

% AV1 gap
\draw[decorate, decoration={brace, amplitude=4pt},
  thick, orange!60!black]
  (4.4, 1.55) -- (4.4, 2.3)
  node[midway, anchor=west, font=\scriptsize, xshift=5pt]
  {$\Delta_{\text{AV1}}$};

% --- Annotation ---
\node[font=\scriptsize\itshape, text=black!60, align=center,
  anchor=north] at (5.0, -0.5)
  {The Shannon rate-distortion bound $R(D)$ is the theoretical minimum.\\
   AV1 narrows the gap $\Delta$ to ${\sim}0.4$ by 2030 (eq.~\ref{eq:coding-gap}).};

\end{tikzpicture}
\caption{Codec operating rates vs.\ the Shannon rate-distortion bound $R(D)$
  at 720p/30fps. Each codec operates above the theoretical minimum; the gap
  $\Delta$ (eq.~\ref{eq:coding-gap}) shrinks from H.264 through VP9 to AV1,
  approaching but never reaching the bound.}
\label{fig:rate-distortion}
\end{figure}

\subsubsection{Codec Entropy and WebCodecs}

The WebCodecs API (W3C Working Draft) exposes per-frame codec access,
allowing measurement and control of the instantaneous encoding rate.
For a video frame $f_t$ at time $t$, the encoder output has empirical
entropy:
\begin{equation}
  \hat{H}(f_t) = -\sum_{s \in \mathcal{S}}
    p(s \mid f_t) \log_2 p(s \mid f_t)
  \label{eq:frame-entropy}
\end{equation}
where $\mathcal{S}$ is the set of symbols emitted by the entropy coder
(CABAC for H.264/AV1, arithmetic coding for VP9). WebCodecs enables
applications to observe $|E(f_t)|$ (the encoded frame size in bits)
directly, which serves as a practical upper bound on
$\hat{H}(f_t)$:
\begin{equation}
  \hat{H}(f_t) \leq |E(f_t)| \leq R_{\text{codec}} / \text{fps}
  \label{eq:entropy-bound}
\end{equation}

\subsection{Effective Information Rate by Time Horizon}

We now combine transport and source coding measures to estimate the
effective information rate for a representative real-time video
session: a single 720p/30fps video call with one audio channel.

\excursus{Excursus: Why 720p?}{Throughout this appendix, 720p/30fps is used as
a fixed reference for cross-era comparison, not as a prediction that 720p
remains the dominant conferencing resolution. The bandwidth surplus quantified
below is precisely what enables migration to 1080p (feasible by ${\sim}$2028
with AV1 hardware encode at today's 720p bit rates) and 4K (feasible by
${\sim}$2030). The 720p baseline makes the efficiency gains legible: the same
task gets cheaper, freeing capacity for higher quality or more streams.}

Define the \emph{net information rate} as:
\begin{equation}
  I_{\text{net}} = R_{\text{codec}} \cdot \eta_{\text{transport}}
    \cdot (1 - \rho_{\text{overhead}})
    \cdot (1 - p_{\text{loss}})
  \label{eq:inet}
\end{equation}

\begin{table}[ht]
\centering
\small
\renewcommand{\arraystretch}{1.4}
\begin{tabularx}{\textwidth}{lXXX}
\toprule
\textbf{Parameter} & \textbf{March 2026} & \textbf{2027--2028} &
  \textbf{2029--2031} \\
\midrule
Median link capacity &
  120~Mbps &
  $\sim$150~Mbps &
  200--300~Mbps \\
Dominant codec (WebRTC) &
  VP9 / H.264 &
  VP9 / AV1 (HW) &
  AV1 (HW default) \\
Codec rate (720p/30fps) &
  1.5--2.5~Mbps &
  1.0--1.8~Mbps &
  0.7--1.2~Mbps \\
$\eta_{\text{transport}}$ &
  0.92 (QUIC) &
  0.93 &
  0.94 \\
$\rho_{\text{overhead}}$ (SRTP+QUIC) &
  0.06 &
  0.05 &
  0.04 \\
Typical loss $p$ &
  0.01--0.03 &
  0.01--0.02 &
  0.005--0.01 \\
$I_{\text{net}}$ (720p video) &
  1.3--2.2~Mbps &
  0.9--1.6~Mbps &
  0.6--1.1~Mbps \\
$I_{\text{net}}$ as \% of $C_{\text{link}}$ &
  1.1--1.8\% &
  0.6--1.1\% &
  0.2--0.6\% \\
\bottomrule
\end{tabularx}
\caption{Effective information rate for a 720p/30fps WebRTC video stream.
  The declining ratio $I_{\text{net}} / C_{\text{link}}$ reflects both
  improved codec efficiency and growing link capacity, freeing bandwidth
  for higher resolutions, additional streams, or concurrent data
  channels.}
\label{tab:inet}
\end{table}

The declining ratio $I_{\text{net}} / C_{\text{link}}$ in
Table~\ref{tab:inet} is the central quantitative observation: by 2030, a
720p video call consumes less than 0.5\% of available link capacity. The
surplus capacity enables qualitative changes --- 4K video conferencing,
multi-stream layouts, or concurrent data channels via WebTransport ---
that are quantitatively infeasible in 2026.

\begin{figure}[ht]
\centering
\begin{tikzpicture}[>=Stealth]

% --- Axes ---
\draw[->, thick] (0, 0) -- (11, 0);
\draw[->, thick] (0, 0) -- (0, 5.5);

% --- Axis title ---
\node[font=\small, rotate=90] at (-0.9, 2.75) {Mbps};

% --- Y-axis labels ---
\foreach \y/\lab in {0/0, 1/50, 2/100, 3/150, 4/200, 5/250} {
  \draw (-0.1, \y) -- (0.1, \y);
  \node[font=\tiny, anchor=east] at (-0.15, \y) {\lab};
}

% --- Group: March 2026 ---
% C_link = 120 Mbps -> y = 2.4
\fill[blue!12] (0.8, 0) rectangle (2.4, 2.4);
\draw[blue!50, thick] (0.8, 0) rectangle (2.4, 2.4);
% I_net = 1.75 Mbps (midpoint) -> y = 0.035
\fill[orange!60] (0.8, 0) rectangle (2.4, 0.07);
\draw[orange!80!black, thick] (0.8, 0) rectangle (2.4, 0.07);
\node[font=\footnotesize, anchor=south] at (1.6, 2.5) {120};
\node[font=\tiny, orange!80!black] at (1.6, -0.4) {1.5\%};
\node[font=\small, anchor=north] at (1.6, -0.7) {2026};

% --- Group: 2027--2028 ---
% C_link = 150 Mbps -> y = 3.0
\fill[blue!12] (4.0, 0) rectangle (5.6, 3.0);
\draw[blue!50, thick] (4.0, 0) rectangle (5.6, 3.0);
% I_net = 1.25 Mbps (midpoint) -> y = 0.025
\fill[orange!60] (4.0, 0) rectangle (5.6, 0.05);
\draw[orange!80!black, thick] (4.0, 0) rectangle (5.6, 0.05);
\node[font=\footnotesize, anchor=south] at (4.8, 3.1) {150};
\node[font=\tiny, orange!80!black] at (4.8, -0.4) {0.8\%};
\node[font=\small, anchor=north] at (4.8, -0.7) {2027--28};

% --- Group: 2029--2031 ---
% C_link = 250 Mbps -> y = 5.0
\fill[blue!12] (7.2, 0) rectangle (8.8, 5.0);
\draw[blue!50, thick] (7.2, 0) rectangle (8.8, 5.0);
% I_net = 0.85 Mbps (midpoint) -> y = 0.017
\fill[orange!60] (7.2, 0) rectangle (8.8, 0.034);
\draw[orange!80!black, thick] (7.2, 0) rectangle (8.8, 0.034);
\node[font=\footnotesize, anchor=south] at (8.0, 5.1) {250};
\node[font=\tiny, orange!80!black] at (8.0, -0.4) {0.3\%};
\node[font=\small, anchor=north] at (8.0, -0.7) {2029--31};

% --- Legend ---
\node[anchor=west, font=\footnotesize] at (0.5, 5.8) {%
  \tikz[baseline=-0.5ex]{\fill[blue!12] (0,0) rectangle (0.3,0.3);
    \draw[blue!50, thick] (0,0) rectangle (0.3,0.3);}
  ~$C_{\text{link}}$ (channel capacity) \qquad
  \tikz[baseline=-0.5ex]{\fill[orange!60] (0,0) rectangle (0.3,0.3);
    \draw[orange!80!black, thick] (0,0) rectangle (0.3,0.3);}
  ~$I_{\text{net}}$ (720p stream)};

% --- Annotation ---
\node[font=\scriptsize\itshape, text=black!60, align=center,
  anchor=north] at (5.5, -1.2)
  {A single 720p call uses a vanishing fraction of Shannon-derived\\
   channel capacity, freeing bandwidth for 4K, multi-stream, and data channels.};

\end{tikzpicture}
\caption{Channel capacity $C_{\text{link}}$ vs.\ effective information rate
  $I_{\text{net}}$ for a single 720p/30fps stream across time horizons. The
  orange bar (stream rate) is barely visible against the blue bar (capacity),
  illustrating the growing surplus as codecs improve and bandwidth increases.}
\label{fig:inet-capacity}
\end{figure}

\subsection{Security Entropy and Key Space}

The DTLS-SRTP security architecture (RFCs~8827, 3711) protects all
WebRTC media. Information-theoretic security requires that the key
entropy exceed the information content of the protected material.

For a DTLS~1.3 handshake with ECDHE on P-256, the key agreement
provides:
\begin{equation}
  H_{\text{key}} = 128 \text{~bits of security}
  \label{eq:key-entropy}
\end{equation}

The SRTP master key is 128 bits (AES-128-GCM), giving an exhaustive
search cost of $2^{128}$ operations. The total information content of a
one-hour 720p video call at 2~Mbps is:
\begin{equation}
  I_{\text{call}} = 2 \times 10^6 \times 3600
    = 7.2 \times 10^9 \text{~bits}
    \approx 2^{32.7} \text{~bits}
  \label{eq:call-content}
\end{equation}

The ratio of key entropy to protected content:
\begin{equation}
  \frac{H_{\text{key}}}{I_{\text{call}}}
    = \frac{128}{32.7} \approx 3.91
  \label{eq:security-margin}
\end{equation}

This ratio remains comfortably above~1 even for much longer sessions.
For a continuous 24-hour 4K AV1 stream at 15~Mbps:
\begin{equation}
  I_{\text{stream}} = 15 \times 10^6 \times 86400
    = 1.296 \times 10^{12} \text{~bits}
    \approx 2^{40.2} \text{~bits}
\end{equation}
yielding $H_{\text{key}} / I_{\text{stream}} \approx 3.18$, still
well above the threshold.

\subsubsection{End-to-End Encryption via Encoded Transforms}

The RTCRtpScriptTransform API (Encoded Transforms) enables
application-layer E2EE within the WebRTC pipeline. From an
information-theoretic perspective, this introduces a second encryption
layer with its own key, creating a \emph{nested channel}:
\begin{equation}
  X \xrightarrow{\;K_{\text{E2EE}}\;}
  E_1(X) \xrightarrow{\;K_{\text{SRTP}}\;}
  E_2(E_1(X))
  \label{eq:nested}
\end{equation}

The mutual information between the ciphertext and plaintext, given
neither key, satisfies:
\begin{equation}
  I\!\left(X;\, E_2(E_1(X))\right) \leq
    \min\!\big(H(K_{\text{E2EE}})^{-1},\;
    H(K_{\text{SRTP}})^{-1}\big)
    \cdot I_{\text{call}}
  \label{eq:mutual-e2ee}
\end{equation}

In practice, with $H(K_{\text{E2EE}}) = H(K_{\text{SRTP}}) = 128$
bits, this quantity is negligible ($\ll 1$ bit for any practical
session length), confirming that the nested encryption provides
information-theoretic confidentiality well beyond computational
security guarantees.

Today (March~2026), Encoded Transforms remain Chrome-only. By
2027--2028 (Section~8.3), cross-browser standardization will make
this nested channel universally available. By 2029--2031, SFrame
for group calls will extend this to $n$-party sessions, where the
mutual information between any non-member and the group plaintext
remains negligible:
\begin{equation}
  I\!\left(X;\, E(X) \;\middle|\; K_i \notin
    \{K_1, \ldots, K_n\}\right) \approx 0
  \label{eq:sframe}
\end{equation}

\subsection{WebTransport: Unreliable Datagrams and Channel Capacity}

WebTransport (RFC~9297) introduces unreliable datagrams to browser-server
communication. For latency-sensitive applications (game state, sensor
telemetry), the relevant measure is not throughput but
\emph{freshness-weighted information rate}. Define the value of a
datagram as exponentially decaying with delivery latency $\delta$:
\begin{equation}
  V(\delta) = e^{-\lambda \delta}
  \label{eq:freshness}
\end{equation}
where $\lambda > 0$ is an application-specific decay rate. The
effective information rate under unreliable delivery is:
\begin{equation}
  I_{\text{WT}} = R_{\text{send}} \cdot (1 - p_{\text{loss}})
    \cdot \mathbb{E}\!\left[e^{-\lambda \delta}\right]
  \label{eq:iwt}
\end{equation}

Under reliable delivery (TCP/WebSocket), a retransmission delays
subsequent data, so:
\begin{equation}
  I_{\text{WS}} = R_{\text{send}}
    \cdot \mathbb{E}\!\left[e^{-\lambda(\delta + p \cdot 2\tau)}\right]
  \label{eq:iws}
\end{equation}

The ratio $I_{\text{WT}} / I_{\text{WS}}$ exceeds 1 when:
\begin{equation}
  (1 - p) \cdot e^{\lambda \cdot p \cdot 2\tau} > 1
  \quad \Longleftrightarrow \quad
  \lambda > \frac{-\ln(1 - p)}{2 p \tau}
  \label{eq:wt-advantage}
\end{equation}

For $p = 0.02$ and $\tau = 25$~ms, this gives $\lambda > 20.2$~s$^{-1}$
(half-life $< 34$~ms).

\excursus{Excursus: Why $p = 0.02$ and $\tau = 25$\,ms?}{These values represent
median fixed-broadband conditions in 2026 (Table~\ref{tab:inet}). On mobile or
congested links ($p \approx 0.05$, $\tau \approx 50$\,ms), the threshold
$\lambda^*$ drops to ${\sim}10$\,s$^{-1}$, broadening WebTransport's advantage
to applications with data half-lives up to ${\sim}70$\,ms. The fixed-broadband
reference is thus conservative: mobile scenarios favor unreliable delivery even
more strongly.}

Any application where data older than
$\sim$30~ms is useless benefits from WebTransport's unreliable
mode over WebSocket --- precisely the regime of interactive gaming, live
cursor tracking, and real-time sensor feeds.

As of March~2026, Safari does not support WebTransport
(Section~5), limiting coverage to $\sim$82\%. By 2027--2028
(Section~8.2), universal support will allow applications to rely on
(\ref{eq:iwt}) without fallback paths.

\begin{figure}[ht]
\centering
\begin{tikzpicture}[>=Stealth]

% --- Axes ---
\draw[->, thick] (0, 0) -- (10.5, 0);
\draw[->, thick] (0, 0) -- (0, 5.5);

% --- Axis titles (positioned to clear tick labels) ---
\node[font=\small] at (5, -0.75)
  {Freshness decay rate $\lambda$ [s$^{-1}$]};
\node[font=\small, rotate=90] at (-1.1, 2.75)
  {$I_{\text{WT}} / I_{\text{WS}}$};

% --- X-axis labels ---
\foreach \x/\lab in {0/0, 1.67/10, 3.33/20, 5.0/30, 6.67/40, 8.33/50, 10/60} {
  \draw (\x, -0.1) -- (\x, 0.1);
  \node[font=\tiny, anchor=north] at (\x, -0.15) {\lab};
}

% --- Y-axis labels (0.96 to 1.06, each 0.02 = 1 unit) ---
\foreach \y/\lab in {0/0.96, 1/0.98, 2/1.00, 3/1.02, 4/1.04, 5/1.06} {
  \draw (-0.1, \y) -- (0.1, \y);
  \node[font=\tiny, anchor=east] at (-0.15, \y) {\lab};
}

% --- Ratio = 1.0 line (y=2 maps to ratio 1.0) ---
\draw[thin, gray, dashed] (0, 2) -- (10, 2);

% --- Shade WebTransport advantage region ---
% lambda* = 20.2 -> x = 3.37
\fill[green!10] (3.37, 2) rectangle (10, 5.3);

% --- I_WT/I_WS curve: 0.98*exp(lambda*0.001) ---
% Y range: 0.96-1.06 (each 0.02 = 1 unit). lambda=0: y=1, lambda*=20.2: y=2, lambda=60: y=4.05
\draw[very thick, orange!80!black]
  (0, 1.0) .. controls (1.5, 1.3) and (2.5, 1.7) ..
  (3.37, 2.0) .. controls (4.5, 2.4) and (6, 2.8) ..
  (8, 3.4) .. controls (9, 3.7) and (9.5, 3.85) ..
  (10, 4.05);

% --- Threshold line ---
\draw[thick, dashed, red!60!black] (3.37, 0) -- (3.37, 5.2);
\node[font=\scriptsize, red!60!black, anchor=south, rotate=90] at (3.17, 3.5)
  {$\lambda^* = 20.2$ s$^{-1}$};

% --- Application regime labels ---
\node[font=\scriptsize\itshape, text=blue!60, anchor=north] at (1.5, 5.3)
  {video conf};
\node[font=\scriptsize\itshape, text=green!50!black, anchor=north] at (7.0, 5.3)
  {gaming / cursor / sensors};

% --- Mark advantage/disadvantage ---
\node[font=\tiny, text=red!50!black, anchor=north east] at (3.2, 1.8)
  {WebSocket wins};
\node[font=\tiny, text=green!50!black, anchor=north west] at (3.5, 4.5)
  {WebTransport wins};

% --- Annotation ---
\node[font=\scriptsize\itshape, text=black!60, align=center,
  anchor=north] at (5.0, -1.0)
  {For $p = 0.02$, $\tau = 25$~ms. Applications where data older than\\
   ${\sim}$30~ms is stale gain from unreliable delivery (eq.~\ref{eq:wt-advantage}).};

\end{tikzpicture}
\caption{Freshness-weighted information rate ratio $I_{\text{WT}} / I_{\text{WS}}$
  as a function of decay rate $\lambda$. Above the threshold $\lambda^*$,
  WebTransport's unreliable datagrams deliver more useful information per
  unit time than WebSocket's reliable delivery, because retransmission
  wastes capacity on stale data.}
\label{fig:freshness}
\end{figure}

\subsection{Projections Summary and Practical Implications}

\begin{table}[ht]
\centering
\small
\renewcommand{\arraystretch}{1.4}
\begin{tabularx}{\textwidth}{lXXX}
\toprule
\textbf{Measure} & \textbf{March 2026} & \textbf{2027--2028} &
  \textbf{2029--2031} \\
\midrule
Channel capacity $C_{\text{link}}$ &
  120~Mbps &
  150~Mbps &
  250~Mbps \\
HOL blocking factor $(1\!-\!p)^{n-2}$ at $n\!=\!8$ &
  0.94 ($p\!=\!0.01$) &
  0.94 &
  0.97 ($p\!=\!0.005$) \\
Source coding gap $\Delta$ &
  VP9: $\sim$0.8 &
  AV1 HW: $\sim$0.5 &
  AV1 opt: $\sim$0.4 \\
$I_{\text{net}}$ per 720p stream &
  1.3--2.2~Mbps &
  0.9--1.6~Mbps &
  0.6--1.1~Mbps \\
Max 720p streams in $C_{\text{link}}$ &
  $\sim$55--90 &
  $\sim$95--165 &
  $\sim$230--415 \\
E2EE availability &
  Chrome only &
  Cross-browser &
  SFrame $n$-party \\
WebTransport coverage &
  82\% &
  $\sim$100\% &
  100\% \\
Freshness threshold $\lambda^*$ &
  20.2~s$^{-1}$ &
  $\sim$18~s$^{-1}$ &
  $\sim$15~s$^{-1}$ \\
\bottomrule
\end{tabularx}
\caption{Information-theoretic measures across time horizons. $\Delta$
  is the coding efficiency gap (\ref{eq:coding-gap}), $I_{\text{net}}$
  is the net information rate (\ref{eq:inet}), and $\lambda^*$ is the
  freshness decay threshold above which WebTransport dominates WebSocket
  (\ref{eq:wt-advantage}).}
\label{tab:projections}
\end{table}

The trajectory is clear: the combination of improving source coding
(AV1 hardware), expanding channel capacity (broadband growth), and
maturing transport protocols (QUIC universal, WebTransport universal)
drives the \emph{information cost per unit of communicated experience}
steadily downward. The remaining question is not whether the
Shannon capacity suffices --- it already does, by orders of magnitude
--- but how efficiently the protocol stack converts that capacity into
delivered, fresh, secure information at the application layer. By
2030, the gap between theoretical capacity and realized throughput will
be the narrowest it has ever been for browser-based real-time
communication.

\subsubsection{What the Theory Means for Applications}

The equations and figures in this appendix are not purely academic. Each
theoretical result maps to a concrete design decision that application
developers face today and over the next five years.

\textbf{Codec selection determines bandwidth cost, not quality.}
The rate-distortion analysis (Section~A.2, Figure~\ref{fig:rate-distortion})
shows that AV1 operates at roughly 46\% of H.264's bit rate for equivalent
perceptual quality at 720p. In practice, this means a WebRTC application
switching from H.264 to AV1 can either halve its bandwidth consumption at
the same quality or substantially increase resolution (e.g., 720p to 1080p)
within the same bandwidth envelope. By 2030, with AV1 hardware encoding
ubiquitous and the coding gap $\Delta$ approaching 0.4, further codec
improvements yield diminishing returns --- the stack will be operating close
enough to the Shannon rate-distortion bound that the next generation of
efficiency gains will come from transport and architecture, not codecs alone.

\textbf{QUIC's advantage is situational but decisive on bad networks.}
The head-of-line blocking analysis (Section~A.1.1,
Figure~\ref{fig:hol-blocking}) shows that QUIC's per-stream throughput
advantage over TCP is modest on good networks (6\% at 1\% loss) but
substantial on degraded ones (34\% at 5\% loss). The practical implication
is that applications targeting mobile users, emerging markets, or congested
enterprise networks --- precisely the environments where real-time
communication is hardest --- benefit disproportionately from QUIC-based
transports. For a video conference with 8 participants on a cellular
connection experiencing 3--5\% packet loss, the difference between QUIC and
TCP is the difference between usable and degraded video.

\textbf{Connection latency is a solved problem for returning users.}
The handshake analysis (Section~A.1.2) shows that QUIC 0-RTT eliminates
connection establishment overhead entirely for resumed connections. For a
WebTransport application that a user opens daily --- a collaboration tool,
a live dashboard, a recurring meeting client --- the time to first useful
byte drops to zero. The 1.125~MB of wasted capacity during a TCP+TLS
handshake on a 120~Mbps link may seem small, but for latency-sensitive
applications where the first 100~ms of perceived responsiveness determines
user experience, the difference is measurable.

\textbf{Bandwidth surplus enables architectural ambition.}
The effective information rate analysis (Section~A.3,
Figure~\ref{fig:inet-capacity}) reveals the most consequential practical
finding: a single 720p video stream consumes a vanishing fraction of
available link capacity (under 0.5\% by 2030). This surplus is not merely
``headroom'' --- it is the resource that makes new application architectures
feasible. A 250~Mbps link can theoretically carry 230--415 simultaneous 720p
streams. In practice, this means applications can afford to run multiple
concurrent media flows (video, screen share, spatial audio), data channels
(collaborative editing, file transfer), and telemetry streams without
contending for bandwidth. The constraint shifts from ``can the network carry
this?''\ to ``can the application manage this complexity?''

\textbf{End-to-end encryption adds negligible overhead.}
The security analysis (Section~A.4) confirms that the nested DTLS-SRTP plus
E2EE architecture imposes no meaningful information-theoretic cost. The key
entropy ($2^{128}$) exceeds the information content of any practical session
by a factor of at least 3. For application developers, this means E2EE can
be treated as a default rather than a premium feature --- there is no
bandwidth, latency, or capacity penalty for encrypting at the application
layer on top of the transport-layer encryption that WebRTC already mandates.
The only constraint is API availability, which moves from Chrome-only (2026)
to universal (2028).

\textbf{WebTransport unlocks real-time data that WebSocket cannot serve.}
The freshness analysis (Section~A.5, Figure~\ref{fig:freshness}) quantifies
when unreliable delivery outperforms reliable delivery: any data with a
useful half-life below approximately 34~ms benefits from WebTransport's
unreliable datagrams. This covers a wide class of interactive applications
--- game state synchronization, live cursor and pointer sharing, real-time
sensor feeds, collaborative whiteboard strokes --- where retransmitting a
stale update is worse than dropping it. Before WebTransport, these
applications had to choose between WebSocket (reliable but stale under loss)
and WebRTC DataChannel (unreliable mode available but requiring the full
peer-to-peer ICE/DTLS machinery). WebTransport provides the unreliable
mode with a simple client-server model over HTTP/3, removing the
architectural overhead.

\textbf{The overall picture: the stack converges on efficiency.}
Taken together, the information-theoretic measures trace a consistent
trajectory. The protocol stack is converging toward the theoretical limits
established by Shannon: codecs approach the rate-distortion bound, transports
eliminate unnecessary blocking and handshake overhead, and encryption adds
negligible cost. By 2030, the practical distance between what the standards
deliver and what information theory permits will be the smallest it has ever
been for browser-based real-time communication. The consequence for
application developers is that the infrastructure --- standards, hardware,
and networks --- will no longer be the binding constraint. The remaining
challenge is purely architectural: how to compose these efficient, universal
primitives into applications that fully exploit the capacity and low latency
they provide.

%
% Requires in the preamble: amsmath, amssymb, tabularx, booktabs

\section{Information-Theoretic View}
\label{app:info-theory}

This appendix applies information-theoretic measures to the real-time
web communication stack described in this document. We quantify channel
capacity, source coding efficiency, and effective information rates
across three time horizons: today (March~2026), the near term
(2027--2028), and the medium term (2029--2031). The analysis draws on
the standards, codec parameters, and infrastructure projections
developed in the main text.

\subsection{Channel Capacity of the Transport Layer}

The Shannon-Hartley theorem gives the theoretical maximum information
rate for a channel with additive white Gaussian noise:
\begin{equation}
  C = B \log_2\!\left(1 + \frac{S}{N}\right) \quad \text{[bits/s]}
  \label{eq:shannon}
\end{equation}
where $B$ is the channel bandwidth in Hz and $S/N$ is the
signal-to-noise ratio.

While the physical-layer capacity of a broadband link is governed by
(\ref{eq:shannon}), the \emph{effective} capacity available to the
application is reduced by protocol overhead, congestion control, and
transport inefficiencies. We define the effective application-layer
throughput as:
\begin{equation}
  R_{\text{eff}} = C_{\text{link}} \cdot \eta_{\text{transport}}
    \cdot (1 - \rho_{\text{overhead}})
  \label{eq:reff}
\end{equation}
where $C_{\text{link}}$ is the link-layer capacity,
$\eta_{\text{transport}} \in (0,1]$ is the transport efficiency factor
(capturing congestion control, retransmission, and multiplexing
behavior), and $\rho_{\text{overhead}} \in [0,1)$ is the fractional
protocol overhead (headers, encryption, framing).

\subsubsection{Head-of-Line Blocking and Multiplexed Capacity}

Consider $n$ independent streams multiplexed over a single connection.
Under TCP (HTTP/2), a single lost packet blocks all $n$ streams. The
expected throughput per stream under loss rate $p$ is bounded by:
\begin{equation}
  R_{\text{TCP}}^{(i)} \leq \frac{R_{\text{eff}}}{n}
    \cdot (1 - p)^{n-1}
  \label{eq:hol-tcp}
\end{equation}
because a loss in \emph{any} of the $n$ streams stalls all others. Under
QUIC (RFC~9000), streams are independent at the transport layer, so:
\begin{equation}
  R_{\text{QUIC}}^{(i)} \leq \frac{R_{\text{eff}}}{n}
    \cdot (1 - p)
  \label{eq:hol-quic}
\end{equation}
The ratio of effective per-stream throughput is:
\begin{equation}
  \frac{R_{\text{QUIC}}^{(i)}}{R_{\text{TCP}}^{(i)}}
    = \frac{1}{(1-p)^{n-2}}
  \label{eq:hol-ratio}
\end{equation}
\excursus{Excursus: Why $n = 8$?}{In a typical multi-participant video
conference, each peer contributes multiple concurrent streams: a camera video
track, an audio track, and often a screen-share or data channel. A four-person
call with two media tracks each produces $n = 8$ multiplexed streams over a
single connection --- the reference value used throughout this section. Larger
meetings (e.g.\ 10 participants with simulcast layers) push $n$ higher,
amplifying the HOL blocking effect further.}

For $n = 8$ streams (a typical multi-track video conference) at
$p = 0.01$ (1\% loss), this gives a QUIC advantage factor of
$1/(0.99)^6 \approx 1.062$, a modest 6.2\% improvement. At $p = 0.05$
(5\% loss, common on mobile or congested links), the factor is
$1/(0.95)^6 \approx 1.34$, a 34\% throughput advantage per stream.

\begin{figure}[ht]
\centering
\begin{tikzpicture}[>=Stealth]

% --- Shaded region (drawn FIRST so curves appear on top) ---
\fill[orange!8] (2, 0) rectangle (10, 5.3);
\node[font=\tiny\itshape, text=black!40, anchor=north east] at (9.8, 5.2)
  {mobile / congested};
\node[font=\tiny\itshape, text=black!40, anchor=north west] at (0.2, 5.2)
  {fixed broadband};

% --- Grid lines ---
\foreach \y in {1, 2, 3, 4, 5} {
  \draw[gray!20] (0, \y) -- (10, \y);
}

% --- Axes ---
\draw[->, thick] (0, 0) -- (10.5, 0);
\draw[->, thick] (0, 0) -- (0, 5.5);

% --- Axis titles (positioned to clear tick labels) ---
\node[font=\small] at (5, -0.85) {Packet loss rate $p$};
\node[font=\small, rotate=90] at (-0.9, 2.75) {QUIC advantage factor};

% --- X-axis labels ---
\foreach \x/\lab in {0/0, 2/0.02, 4/0.04, 6/0.06, 8/0.08, 10/0.10} {
  \draw (\x, -0.1) -- (\x, 0.1);
  \node[font=\tiny, anchor=north] at (\x, -0.15) {\lab};
}

% --- Y-axis labels ---
\foreach \y/\lab in {0/1.0, 1/1.1, 2/1.2, 3/1.3, 4/1.4, 5/1.5} {
  \draw (-0.1, \y) -- (0.1, \y);
  \node[font=\tiny, anchor=east] at (-0.15, \y) {\lab};
}

% --- n=4 curve: 1/(1-p)^2 ---
\draw[very thick, blue!70]
  (0, 0) .. controls (2.5, 0.5) and (5, 1.1) ..
  (7.5, 2.0) .. controls (8.5, 2.4) and (9.5, 2.9) ..
  (10, 3.1);
\node[font=\scriptsize, blue!70, anchor=south west] at (9.0, 3.0)
  {$n = 4$};

% --- n=8 curve: 1/(1-p)^6 ---
\draw[very thick, orange!80!black]
  (0, 0) .. controls (1.5, 0.5) and (3, 1.7) ..
  (4, 2.5) .. controls (5, 3.5) and (5.5, 4.0) ..
  (6, 4.7);
\node[font=\scriptsize, orange!80!black, anchor=south] at (5.5, 4.8)
  {$n = 8$};

% --- Reference points ---
\fill[orange!80!black] (1.0, 0.62) circle (2.5pt);
\node[font=\tiny, orange!80!black, anchor=south west] at (1.1, 0.62)
  {6.2\%};

\fill[orange!80!black] (5.0, 3.7) circle (2.5pt);
\node[font=\tiny, orange!80!black, anchor=west] at (5.1, 3.7)
  {34\%};

\fill[blue!70] (1.0, 0.20) circle (2.5pt);

% --- Annotation ---
\node[font=\scriptsize\itshape, text=black!60, align=center,
  anchor=north] at (5.0, -1.0)
  {QUIC recovers per-stream throughput closer to $C/n$ (Shannon fair share)\\
   by eliminating head-of-line blocking (eq.~\ref{eq:hol-ratio}).};

\end{tikzpicture}
\caption{QUIC advantage factor $1/(1-p)^{n-2}$ over TCP for multiplexed
  streams. At 5\% loss with $n=8$ streams, QUIC delivers 34\% more
  per-stream throughput by maintaining independence at the transport layer.}
\label{fig:hol-blocking}
\end{figure}

\subsubsection{Connection Establishment Latency}

Let $\tau$ denote the one-way latency (half the round-trip time). The
information-carrying capacity during connection establishment is zero;
the handshake represents pure overhead. The time to first useful byte
is:
\begin{align}
  T_{\text{TCP+TLS\,1.3}} &= 2\tau_{\text{TCP}} + \tau_{\text{TLS}}
    = 3\tau \quad \text{(1-RTT TLS~1.3)} \label{eq:tcp-setup} \\
  T_{\text{QUIC}} &= \tau \quad \text{(1-RTT)}
    \label{eq:quic-setup} \\
  T_{\text{QUIC,\,0-RTT}} &= 0 \quad \text{(resumption)}
    \label{eq:quic-0rtt}
\end{align}
The QUIC 0-RTT case (\ref{eq:quic-0rtt}) eliminates the handshake
entirely for resumed connections. The information lost during
handshake dead time, relative to steady-state capacity, is:
\begin{equation}
  I_{\text{lost}} = C_{\text{link}} \cdot T_{\text{setup}}
  \label{eq:ilost}
\end{equation}
For a 120~Mbps link with $\tau = 25$~ms (typical fixed broadband in
2026), the TCP+TLS~1.3 handshake wastes $120 \times 10^6 \times 0.075
= 9 \times 10^6$ bits (1.125~MB), while QUIC~1-RTT wastes only
$3 \times 10^6$ bits (375~KB).

\subsection{Source Coding Efficiency}

The rate-distortion function $R(D)$ defines the minimum bit rate
required to represent a source at distortion level $D$. For a
Gaussian source with variance $\sigma^2$ under mean squared error
distortion:
\begin{equation}
  R(D) = \frac{1}{2} \log_2 \frac{\sigma^2}{D}
    \quad \text{[bits/sample]}
  \label{eq:rate-distortion}
\end{equation}

\excursus{Excursus: Why Gaussian source?}{The Gaussian memoryless model is used
because it yields the highest (worst-case) $R(D)$ among all sources of equal
variance. Any codec that approaches the Gaussian bound on real video is
therefore doing at least as well as the bound predicts --- the true $R(D)$ for
spatially and temporally correlated video lies below the Gaussian curve, making
it a conservative benchmark.}

No practical codec achieves $R(D)$;\footnote{The $R(D)$ curve in
Figure~\ref{fig:rate-distortion} is the Gaussian memoryless source bound
(equation~\ref{eq:rate-distortion}), which is an upper bound on the true
$R(D)$ for non-Gaussian sources such as video.  Additionally, $R(D)$ assumes
zero excess-distortion probability, whereas practical codecs permit a
frame-level excess-distortion probability on the order of $10^{-6}$ to
$10^{-8}$.  Both factors allow codec operating points measured on real content
to approach or marginally cross the Gaussian bound.} real codecs operate at
some rate $R_{\text{codec}} > R(D)$. We define the coding efficiency gap as:
\begin{equation}
  \Delta = \frac{R_{\text{codec}} - R(D)}{R(D)}
  \label{eq:coding-gap}
\end{equation}

The codecs in the WebRTC stack can be ordered by their empirical
proximity to $R(D)$ for typical video content at conversational
resolutions (720p, 30~fps). Using published Bj\o ntegaard Delta (BD)
rate measurements:
\begin{equation}
  R_{\text{H.264}} > R_{\text{VP9}} > R_{\text{AV1}}
    > R(D)
  \label{eq:codec-ordering}
\end{equation}

Specifically, the empirical rate savings at equivalent PSNR are
approximately:
\begin{align}
  R_{\text{VP9}}  &\approx 0.65 \cdot R_{\text{H.264}}
    \quad (\sim\!35\% \text{ savings}) \label{eq:vp9-savings} \\
  R_{\text{AV1}}  &\approx 0.70 \cdot R_{\text{VP9}}
    \approx 0.46 \cdot R_{\text{H.264}}
    \quad (\sim\!30\% \text{ over VP9})
  \label{eq:av1-savings}
\end{align}

\begin{figure}[ht]
\centering
\begin{tikzpicture}[>=Stealth]

% --- Axes ---
\draw[->, thick] (0, 0) -- (10.5, 0);
\draw[->, thick] (0, 0) -- (0, 6.5);

% --- Axis titles ---
\node[font=\small] at (5, -0.85) {Distortion $D$};
\node[font=\small, rotate=90] at (-0.9, 3.25) {Rate [Mbps]};

% --- Y-axis rate labels ---
\foreach \y/\lab in {1.0/0.5, 2.0/1.0, 3.0/1.5, 4.0/2.0, 5.0/2.5, 6.0/3.0} {
  \draw (-0.1, \y) -- (0.1, \y);
  \node[font=\tiny, anchor=east] at (-0.15, \y) {\lab};
}

% --- R(D) Shannon bound curve ---
\draw[very thick, green!60!black]
  (1.0, 5.8) .. controls (2.0, 3.5) and (3.5, 2.0) ..
  (5.0, 1.3) .. controls (6.5, 0.8) and (8.0, 0.55) ..
  (9.5, 0.4);
\node[font=\scriptsize, green!60!black, anchor=north west] at (8.5, 0.55)
  {$R(D)$ bound};

% --- Reference distortion line (720p/30fps operating point) ---
\draw[thin, gray, dashed] (4.0, 0) -- (4.0, 6.2);
\node[font=\tiny, gray, anchor=south] at (4.0, 6.2) {720p/30fps};

% --- R(D) value at reference distortion ---
\fill[green!60!black] (4.0, 1.55) circle (2pt);

% --- Codec operating points at reference distortion ---
% H.264: ~2.5 Mbps
\draw[thick, dashed, red!70!black] (0.3, 4.8) -- (9.5, 4.8);
\fill[red!70!black] (4.0, 4.8) circle (3pt);
\node[font=\scriptsize, red!70!black, anchor=west] at (9.6, 4.8) {H.264};

% VP9: ~1.6 Mbps (0.65 * H.264)
\draw[thick, dashed, blue!70] (0.3, 3.3) -- (9.5, 3.3);
\fill[blue!70] (4.0, 3.3) circle (3pt);
\node[font=\scriptsize, blue!70, anchor=west] at (9.6, 3.3) {VP9};

% AV1: ~1.15 Mbps (0.46 * H.264)
\draw[thick, dashed, orange!80!black] (0.3, 2.3) -- (9.5, 2.3);
\fill[orange!80!black] (4.0, 2.3) circle (3pt);
\node[font=\scriptsize, orange!80!black, anchor=west] at (9.6, 2.3) {AV1};

% --- Gap annotations (braces) ---
% H.264 gap
\draw[decorate, decoration={brace, amplitude=4pt, mirror},
  thick, red!50!black]
  (3.6, 1.55) -- (3.6, 4.8)
  node[midway, anchor=east, font=\scriptsize, xshift=-5pt]
  {$\Delta_{\text{H.264}}$};

% AV1 gap
\draw[decorate, decoration={brace, amplitude=4pt},
  thick, orange!60!black]
  (4.4, 1.55) -- (4.4, 2.3)
  node[midway, anchor=west, font=\scriptsize, xshift=5pt]
  {$\Delta_{\text{AV1}}$};

% --- Annotation ---
\node[font=\scriptsize\itshape, text=black!60, align=center,
  anchor=north] at (5.0, -0.5)
  {The Shannon rate-distortion bound $R(D)$ is the theoretical minimum.\\
   AV1 narrows the gap $\Delta$ to ${\sim}0.4$ by 2030 (eq.~\ref{eq:coding-gap}).};

\end{tikzpicture}
\caption{Codec operating rates vs.\ the Shannon rate-distortion bound $R(D)$
  at 720p/30fps. Each codec operates above the theoretical minimum; the gap
  $\Delta$ (eq.~\ref{eq:coding-gap}) shrinks from H.264 through VP9 to AV1,
  approaching but never reaching the bound.}
\label{fig:rate-distortion}
\end{figure}

\subsubsection{Codec Entropy and WebCodecs}

The WebCodecs API (W3C Working Draft) exposes per-frame codec access,
allowing measurement and control of the instantaneous encoding rate.
For a video frame $f_t$ at time $t$, the encoder output has empirical
entropy:
\begin{equation}
  \hat{H}(f_t) = -\sum_{s \in \mathcal{S}}
    p(s \mid f_t) \log_2 p(s \mid f_t)
  \label{eq:frame-entropy}
\end{equation}
where $\mathcal{S}$ is the set of symbols emitted by the entropy coder
(CABAC for H.264/AV1, arithmetic coding for VP9). WebCodecs enables
applications to observe $|E(f_t)|$ (the encoded frame size in bits)
directly, which serves as a practical upper bound on
$\hat{H}(f_t)$:
\begin{equation}
  \hat{H}(f_t) \leq |E(f_t)| \leq R_{\text{codec}} / \text{fps}
  \label{eq:entropy-bound}
\end{equation}

\subsection{Effective Information Rate by Time Horizon}

We now combine transport and source coding measures to estimate the
effective information rate for a representative real-time video
session: a single 720p/30fps video call with one audio channel.

\excursus{Excursus: Why 720p?}{Throughout this appendix, 720p/30fps is used as
a fixed reference for cross-era comparison, not as a prediction that 720p
remains the dominant conferencing resolution. The bandwidth surplus quantified
below is precisely what enables migration to 1080p (feasible by ${\sim}$2028
with AV1 hardware encode at today's 720p bit rates) and 4K (feasible by
${\sim}$2030). The 720p baseline makes the efficiency gains legible: the same
task gets cheaper, freeing capacity for higher quality or more streams.}

Define the \emph{net information rate} as:
\begin{equation}
  I_{\text{net}} = R_{\text{codec}} \cdot \eta_{\text{transport}}
    \cdot (1 - \rho_{\text{overhead}})
    \cdot (1 - p_{\text{loss}})
  \label{eq:inet}
\end{equation}

\begin{table}[ht]
\centering
\small
\renewcommand{\arraystretch}{1.4}
\begin{tabularx}{\textwidth}{lXXX}
\toprule
\textbf{Parameter} & \textbf{March 2026} & \textbf{2027--2028} &
  \textbf{2029--2031} \\
\midrule
Median link capacity &
  120~Mbps &
  $\sim$150~Mbps &
  200--300~Mbps \\
Dominant codec (WebRTC) &
  VP9 / H.264 &
  VP9 / AV1 (HW) &
  AV1 (HW default) \\
Codec rate (720p/30fps) &
  1.5--2.5~Mbps &
  1.0--1.8~Mbps &
  0.7--1.2~Mbps \\
$\eta_{\text{transport}}$ &
  0.92 (QUIC) &
  0.93 &
  0.94 \\
$\rho_{\text{overhead}}$ (SRTP+QUIC) &
  0.06 &
  0.05 &
  0.04 \\
Typical loss $p$ &
  0.01--0.03 &
  0.01--0.02 &
  0.005--0.01 \\
$I_{\text{net}}$ (720p video) &
  1.3--2.2~Mbps &
  0.9--1.6~Mbps &
  0.6--1.1~Mbps \\
$I_{\text{net}}$ as \% of $C_{\text{link}}$ &
  1.1--1.8\% &
  0.6--1.1\% &
  0.2--0.6\% \\
\bottomrule
\end{tabularx}
\caption{Effective information rate for a 720p/30fps WebRTC video stream.
  The declining ratio $I_{\text{net}} / C_{\text{link}}$ reflects both
  improved codec efficiency and growing link capacity, freeing bandwidth
  for higher resolutions, additional streams, or concurrent data
  channels.}
\label{tab:inet}
\end{table}

The declining ratio $I_{\text{net}} / C_{\text{link}}$ in
Table~\ref{tab:inet} is the central quantitative observation: by 2030, a
720p video call consumes less than 0.5\% of available link capacity. The
surplus capacity enables qualitative changes --- 4K video conferencing,
multi-stream layouts, or concurrent data channels via WebTransport ---
that are quantitatively infeasible in 2026.

\begin{figure}[ht]
\centering
\begin{tikzpicture}[>=Stealth]

% --- Axes ---
\draw[->, thick] (0, 0) -- (11, 0);
\draw[->, thick] (0, 0) -- (0, 5.5);

% --- Axis title ---
\node[font=\small, rotate=90] at (-0.9, 2.75) {Mbps};

% --- Y-axis labels ---
\foreach \y/\lab in {0/0, 1/50, 2/100, 3/150, 4/200, 5/250} {
  \draw (-0.1, \y) -- (0.1, \y);
  \node[font=\tiny, anchor=east] at (-0.15, \y) {\lab};
}

% --- Group: March 2026 ---
% C_link = 120 Mbps -> y = 2.4
\fill[blue!12] (0.8, 0) rectangle (2.4, 2.4);
\draw[blue!50, thick] (0.8, 0) rectangle (2.4, 2.4);
% I_net = 1.75 Mbps (midpoint) -> y = 0.035
\fill[orange!60] (0.8, 0) rectangle (2.4, 0.07);
\draw[orange!80!black, thick] (0.8, 0) rectangle (2.4, 0.07);
\node[font=\footnotesize, anchor=south] at (1.6, 2.5) {120};
\node[font=\tiny, orange!80!black] at (1.6, -0.4) {1.5\%};
\node[font=\small, anchor=north] at (1.6, -0.7) {2026};

% --- Group: 2027--2028 ---
% C_link = 150 Mbps -> y = 3.0
\fill[blue!12] (4.0, 0) rectangle (5.6, 3.0);
\draw[blue!50, thick] (4.0, 0) rectangle (5.6, 3.0);
% I_net = 1.25 Mbps (midpoint) -> y = 0.025
\fill[orange!60] (4.0, 0) rectangle (5.6, 0.05);
\draw[orange!80!black, thick] (4.0, 0) rectangle (5.6, 0.05);
\node[font=\footnotesize, anchor=south] at (4.8, 3.1) {150};
\node[font=\tiny, orange!80!black] at (4.8, -0.4) {0.8\%};
\node[font=\small, anchor=north] at (4.8, -0.7) {2027--28};

% --- Group: 2029--2031 ---
% C_link = 250 Mbps -> y = 5.0
\fill[blue!12] (7.2, 0) rectangle (8.8, 5.0);
\draw[blue!50, thick] (7.2, 0) rectangle (8.8, 5.0);
% I_net = 0.85 Mbps (midpoint) -> y = 0.017
\fill[orange!60] (7.2, 0) rectangle (8.8, 0.034);
\draw[orange!80!black, thick] (7.2, 0) rectangle (8.8, 0.034);
\node[font=\footnotesize, anchor=south] at (8.0, 5.1) {250};
\node[font=\tiny, orange!80!black] at (8.0, -0.4) {0.3\%};
\node[font=\small, anchor=north] at (8.0, -0.7) {2029--31};

% --- Legend ---
\node[anchor=west, font=\footnotesize] at (0.5, 5.8) {%
  \tikz[baseline=-0.5ex]{\fill[blue!12] (0,0) rectangle (0.3,0.3);
    \draw[blue!50, thick] (0,0) rectangle (0.3,0.3);}
  ~$C_{\text{link}}$ (channel capacity) \qquad
  \tikz[baseline=-0.5ex]{\fill[orange!60] (0,0) rectangle (0.3,0.3);
    \draw[orange!80!black, thick] (0,0) rectangle (0.3,0.3);}
  ~$I_{\text{net}}$ (720p stream)};

% --- Annotation ---
\node[font=\scriptsize\itshape, text=black!60, align=center,
  anchor=north] at (5.5, -1.2)
  {A single 720p call uses a vanishing fraction of Shannon-derived\\
   channel capacity, freeing bandwidth for 4K, multi-stream, and data channels.};

\end{tikzpicture}
\caption{Channel capacity $C_{\text{link}}$ vs.\ effective information rate
  $I_{\text{net}}$ for a single 720p/30fps stream across time horizons. The
  orange bar (stream rate) is barely visible against the blue bar (capacity),
  illustrating the growing surplus as codecs improve and bandwidth increases.}
\label{fig:inet-capacity}
\end{figure}

\subsection{Security Entropy and Key Space}

The DTLS-SRTP security architecture (RFCs~8827, 3711) protects all
WebRTC media. Information-theoretic security requires that the key
entropy exceed the information content of the protected material.

For a DTLS~1.3 handshake with ECDHE on P-256, the key agreement
provides:
\begin{equation}
  H_{\text{key}} = 128 \text{~bits of security}
  \label{eq:key-entropy}
\end{equation}

The SRTP master key is 128 bits (AES-128-GCM), giving an exhaustive
search cost of $2^{128}$ operations. The total information content of a
one-hour 720p video call at 2~Mbps is:
\begin{equation}
  I_{\text{call}} = 2 \times 10^6 \times 3600
    = 7.2 \times 10^9 \text{~bits}
    \approx 2^{32.7} \text{~bits}
  \label{eq:call-content}
\end{equation}

The ratio of key entropy to protected content:
\begin{equation}
  \frac{H_{\text{key}}}{I_{\text{call}}}
    = \frac{128}{32.7} \approx 3.91
  \label{eq:security-margin}
\end{equation}

This ratio remains comfortably above~1 even for much longer sessions.
For a continuous 24-hour 4K AV1 stream at 15~Mbps:
\begin{equation}
  I_{\text{stream}} = 15 \times 10^6 \times 86400
    = 1.296 \times 10^{12} \text{~bits}
    \approx 2^{40.2} \text{~bits}
\end{equation}
yielding $H_{\text{key}} / I_{\text{stream}} \approx 3.18$, still
well above the threshold.

\subsubsection{End-to-End Encryption via Encoded Transforms}

The RTCRtpScriptTransform API (Encoded Transforms) enables
application-layer E2EE within the WebRTC pipeline. From an
information-theoretic perspective, this introduces a second encryption
layer with its own key, creating a \emph{nested channel}:
\begin{equation}
  X \xrightarrow{\;K_{\text{E2EE}}\;}
  E_1(X) \xrightarrow{\;K_{\text{SRTP}}\;}
  E_2(E_1(X))
  \label{eq:nested}
\end{equation}

The mutual information between the ciphertext and plaintext, given
neither key, satisfies:
\begin{equation}
  I\!\left(X;\, E_2(E_1(X))\right) \leq
    \min\!\big(H(K_{\text{E2EE}})^{-1},\;
    H(K_{\text{SRTP}})^{-1}\big)
    \cdot I_{\text{call}}
  \label{eq:mutual-e2ee}
\end{equation}

In practice, with $H(K_{\text{E2EE}}) = H(K_{\text{SRTP}}) = 128$
bits, this quantity is negligible ($\ll 1$ bit for any practical
session length), confirming that the nested encryption provides
information-theoretic confidentiality well beyond computational
security guarantees.

Today (March~2026), Encoded Transforms remain Chrome-only. By
2027--2028 (Section~8.3), cross-browser standardization will make
this nested channel universally available. By 2029--2031, SFrame
for group calls will extend this to $n$-party sessions, where the
mutual information between any non-member and the group plaintext
remains negligible:
\begin{equation}
  I\!\left(X;\, E(X) \;\middle|\; K_i \notin
    \{K_1, \ldots, K_n\}\right) \approx 0
  \label{eq:sframe}
\end{equation}

\subsection{WebTransport: Unreliable Datagrams and Channel Capacity}

WebTransport (RFC~9297) introduces unreliable datagrams to browser-server
communication. For latency-sensitive applications (game state, sensor
telemetry), the relevant measure is not throughput but
\emph{freshness-weighted information rate}. Define the value of a
datagram as exponentially decaying with delivery latency $\delta$:
\begin{equation}
  V(\delta) = e^{-\lambda \delta}
  \label{eq:freshness}
\end{equation}
where $\lambda > 0$ is an application-specific decay rate. The
effective information rate under unreliable delivery is:
\begin{equation}
  I_{\text{WT}} = R_{\text{send}} \cdot (1 - p_{\text{loss}})
    \cdot \mathbb{E}\!\left[e^{-\lambda \delta}\right]
  \label{eq:iwt}
\end{equation}

Under reliable delivery (TCP/WebSocket), a retransmission delays
subsequent data, so:
\begin{equation}
  I_{\text{WS}} = R_{\text{send}}
    \cdot \mathbb{E}\!\left[e^{-\lambda(\delta + p \cdot 2\tau)}\right]
  \label{eq:iws}
\end{equation}

The ratio $I_{\text{WT}} / I_{\text{WS}}$ exceeds 1 when:
\begin{equation}
  (1 - p) \cdot e^{\lambda \cdot p \cdot 2\tau} > 1
  \quad \Longleftrightarrow \quad
  \lambda > \frac{-\ln(1 - p)}{2 p \tau}
  \label{eq:wt-advantage}
\end{equation}

For $p = 0.02$ and $\tau = 25$~ms, this gives $\lambda > 20.2$~s$^{-1}$
(half-life $< 34$~ms).

\excursus{Excursus: Why $p = 0.02$ and $\tau = 25$\,ms?}{These values represent
median fixed-broadband conditions in 2026 (Table~\ref{tab:inet}). On mobile or
congested links ($p \approx 0.05$, $\tau \approx 50$\,ms), the threshold
$\lambda^*$ drops to ${\sim}10$\,s$^{-1}$, broadening WebTransport's advantage
to applications with data half-lives up to ${\sim}70$\,ms. The fixed-broadband
reference is thus conservative: mobile scenarios favor unreliable delivery even
more strongly.}

Any application where data older than
$\sim$30~ms is useless benefits from WebTransport's unreliable
mode over WebSocket --- precisely the regime of interactive gaming, live
cursor tracking, and real-time sensor feeds.

As of March~2026, Safari does not support WebTransport
(Section~5), limiting coverage to $\sim$82\%. By 2027--2028
(Section~8.2), universal support will allow applications to rely on
(\ref{eq:iwt}) without fallback paths.

\begin{figure}[ht]
\centering
\begin{tikzpicture}[>=Stealth]

% --- Axes ---
\draw[->, thick] (0, 0) -- (10.5, 0);
\draw[->, thick] (0, 0) -- (0, 5.5);

% --- Axis titles (positioned to clear tick labels) ---
\node[font=\small] at (5, -0.75)
  {Freshness decay rate $\lambda$ [s$^{-1}$]};
\node[font=\small, rotate=90] at (-1.1, 2.75)
  {$I_{\text{WT}} / I_{\text{WS}}$};

% --- X-axis labels ---
\foreach \x/\lab in {0/0, 1.67/10, 3.33/20, 5.0/30, 6.67/40, 8.33/50, 10/60} {
  \draw (\x, -0.1) -- (\x, 0.1);
  \node[font=\tiny, anchor=north] at (\x, -0.15) {\lab};
}

% --- Y-axis labels (0.96 to 1.06, each 0.02 = 1 unit) ---
\foreach \y/\lab in {0/0.96, 1/0.98, 2/1.00, 3/1.02, 4/1.04, 5/1.06} {
  \draw (-0.1, \y) -- (0.1, \y);
  \node[font=\tiny, anchor=east] at (-0.15, \y) {\lab};
}

% --- Ratio = 1.0 line (y=2 maps to ratio 1.0) ---
\draw[thin, gray, dashed] (0, 2) -- (10, 2);

% --- Shade WebTransport advantage region ---
% lambda* = 20.2 -> x = 3.37
\fill[green!10] (3.37, 2) rectangle (10, 5.3);

% --- I_WT/I_WS curve: 0.98*exp(lambda*0.001) ---
% Y range: 0.96-1.06 (each 0.02 = 1 unit). lambda=0: y=1, lambda*=20.2: y=2, lambda=60: y=4.05
\draw[very thick, orange!80!black]
  (0, 1.0) .. controls (1.5, 1.3) and (2.5, 1.7) ..
  (3.37, 2.0) .. controls (4.5, 2.4) and (6, 2.8) ..
  (8, 3.4) .. controls (9, 3.7) and (9.5, 3.85) ..
  (10, 4.05);

% --- Threshold line ---
\draw[thick, dashed, red!60!black] (3.37, 0) -- (3.37, 5.2);
\node[font=\scriptsize, red!60!black, anchor=south, rotate=90] at (3.17, 3.5)
  {$\lambda^* = 20.2$ s$^{-1}$};

% --- Application regime labels ---
\node[font=\scriptsize\itshape, text=blue!60, anchor=north] at (1.5, 5.3)
  {video conf};
\node[font=\scriptsize\itshape, text=green!50!black, anchor=north] at (7.0, 5.3)
  {gaming / cursor / sensors};

% --- Mark advantage/disadvantage ---
\node[font=\tiny, text=red!50!black, anchor=north east] at (3.2, 1.8)
  {WebSocket wins};
\node[font=\tiny, text=green!50!black, anchor=north west] at (3.5, 4.5)
  {WebTransport wins};

% --- Annotation ---
\node[font=\scriptsize\itshape, text=black!60, align=center,
  anchor=north] at (5.0, -1.0)
  {For $p = 0.02$, $\tau = 25$~ms. Applications where data older than\\
   ${\sim}$30~ms is stale gain from unreliable delivery (eq.~\ref{eq:wt-advantage}).};

\end{tikzpicture}
\caption{Freshness-weighted information rate ratio $I_{\text{WT}} / I_{\text{WS}}$
  as a function of decay rate $\lambda$. Above the threshold $\lambda^*$,
  WebTransport's unreliable datagrams deliver more useful information per
  unit time than WebSocket's reliable delivery, because retransmission
  wastes capacity on stale data.}
\label{fig:freshness}
\end{figure}

\subsection{Projections Summary and Practical Implications}

\begin{table}[ht]
\centering
\small
\renewcommand{\arraystretch}{1.4}
\begin{tabularx}{\textwidth}{lXXX}
\toprule
\textbf{Measure} & \textbf{March 2026} & \textbf{2027--2028} &
  \textbf{2029--2031} \\
\midrule
Channel capacity $C_{\text{link}}$ &
  120~Mbps &
  150~Mbps &
  250~Mbps \\
HOL blocking factor $(1\!-\!p)^{n-2}$ at $n\!=\!8$ &
  0.94 ($p\!=\!0.01$) &
  0.94 &
  0.97 ($p\!=\!0.005$) \\
Source coding gap $\Delta$ &
  VP9: $\sim$0.8 &
  AV1 HW: $\sim$0.5 &
  AV1 opt: $\sim$0.4 \\
$I_{\text{net}}$ per 720p stream &
  1.3--2.2~Mbps &
  0.9--1.6~Mbps &
  0.6--1.1~Mbps \\
Max 720p streams in $C_{\text{link}}$ &
  $\sim$55--90 &
  $\sim$95--165 &
  $\sim$230--415 \\
E2EE availability &
  Chrome only &
  Cross-browser &
  SFrame $n$-party \\
WebTransport coverage &
  82\% &
  $\sim$100\% &
  100\% \\
Freshness threshold $\lambda^*$ &
  20.2~s$^{-1}$ &
  $\sim$18~s$^{-1}$ &
  $\sim$15~s$^{-1}$ \\
\bottomrule
\end{tabularx}
\caption{Information-theoretic measures across time horizons. $\Delta$
  is the coding efficiency gap (\ref{eq:coding-gap}), $I_{\text{net}}$
  is the net information rate (\ref{eq:inet}), and $\lambda^*$ is the
  freshness decay threshold above which WebTransport dominates WebSocket
  (\ref{eq:wt-advantage}).}
\label{tab:projections}
\end{table}

The trajectory is clear: the combination of improving source coding
(AV1 hardware), expanding channel capacity (broadband growth), and
maturing transport protocols (QUIC universal, WebTransport universal)
drives the \emph{information cost per unit of communicated experience}
steadily downward. The remaining question is not whether the
Shannon capacity suffices --- it already does, by orders of magnitude
--- but how efficiently the protocol stack converts that capacity into
delivered, fresh, secure information at the application layer. By
2030, the gap between theoretical capacity and realized throughput will
be the narrowest it has ever been for browser-based real-time
communication.

\subsubsection{What the Theory Means for Applications}

The equations and figures in this appendix are not purely academic. Each
theoretical result maps to a concrete design decision that application
developers face today and over the next five years.

\textbf{Codec selection determines bandwidth cost, not quality.}
The rate-distortion analysis (Section~A.2, Figure~\ref{fig:rate-distortion})
shows that AV1 operates at roughly 46\% of H.264's bit rate for equivalent
perceptual quality at 720p. In practice, this means a WebRTC application
switching from H.264 to AV1 can either halve its bandwidth consumption at
the same quality or substantially increase resolution (e.g., 720p to 1080p)
within the same bandwidth envelope. By 2030, with AV1 hardware encoding
ubiquitous and the coding gap $\Delta$ approaching 0.4, further codec
improvements yield diminishing returns --- the stack will be operating close
enough to the Shannon rate-distortion bound that the next generation of
efficiency gains will come from transport and architecture, not codecs alone.

\textbf{QUIC's advantage is situational but decisive on bad networks.}
The head-of-line blocking analysis (Section~A.1.1,
Figure~\ref{fig:hol-blocking}) shows that QUIC's per-stream throughput
advantage over TCP is modest on good networks (6\% at 1\% loss) but
substantial on degraded ones (34\% at 5\% loss). The practical implication
is that applications targeting mobile users, emerging markets, or congested
enterprise networks --- precisely the environments where real-time
communication is hardest --- benefit disproportionately from QUIC-based
transports. For a video conference with 8 participants on a cellular
connection experiencing 3--5\% packet loss, the difference between QUIC and
TCP is the difference between usable and degraded video.

\textbf{Connection latency is a solved problem for returning users.}
The handshake analysis (Section~A.1.2) shows that QUIC 0-RTT eliminates
connection establishment overhead entirely for resumed connections. For a
WebTransport application that a user opens daily --- a collaboration tool,
a live dashboard, a recurring meeting client --- the time to first useful
byte drops to zero. The 1.125~MB of wasted capacity during a TCP+TLS
handshake on a 120~Mbps link may seem small, but for latency-sensitive
applications where the first 100~ms of perceived responsiveness determines
user experience, the difference is measurable.

\textbf{Bandwidth surplus enables architectural ambition.}
The effective information rate analysis (Section~A.3,
Figure~\ref{fig:inet-capacity}) reveals the most consequential practical
finding: a single 720p video stream consumes a vanishing fraction of
available link capacity (under 0.5\% by 2030). This surplus is not merely
``headroom'' --- it is the resource that makes new application architectures
feasible. A 250~Mbps link can theoretically carry 230--415 simultaneous 720p
streams. In practice, this means applications can afford to run multiple
concurrent media flows (video, screen share, spatial audio), data channels
(collaborative editing, file transfer), and telemetry streams without
contending for bandwidth. The constraint shifts from ``can the network carry
this?''\ to ``can the application manage this complexity?''

\textbf{End-to-end encryption adds negligible overhead.}
The security analysis (Section~A.4) confirms that the nested DTLS-SRTP plus
E2EE architecture imposes no meaningful information-theoretic cost. The key
entropy ($2^{128}$) exceeds the information content of any practical session
by a factor of at least 3. For application developers, this means E2EE can
be treated as a default rather than a premium feature --- there is no
bandwidth, latency, or capacity penalty for encrypting at the application
layer on top of the transport-layer encryption that WebRTC already mandates.
The only constraint is API availability, which moves from Chrome-only (2026)
to universal (2028).

\textbf{WebTransport unlocks real-time data that WebSocket cannot serve.}
The freshness analysis (Section~A.5, Figure~\ref{fig:freshness}) quantifies
when unreliable delivery outperforms reliable delivery: any data with a
useful half-life below approximately 34~ms benefits from WebTransport's
unreliable datagrams. This covers a wide class of interactive applications
--- game state synchronization, live cursor and pointer sharing, real-time
sensor feeds, collaborative whiteboard strokes --- where retransmitting a
stale update is worse than dropping it. Before WebTransport, these
applications had to choose between WebSocket (reliable but stale under loss)
and WebRTC DataChannel (unreliable mode available but requiring the full
peer-to-peer ICE/DTLS machinery). WebTransport provides the unreliable
mode with a simple client-server model over HTTP/3, removing the
architectural overhead.

\textbf{The overall picture: the stack converges on efficiency.}
Taken together, the information-theoretic measures trace a consistent
trajectory. The protocol stack is converging toward the theoretical limits
established by Shannon: codecs approach the rate-distortion bound, transports
eliminate unnecessary blocking and handshake overhead, and encryption adds
negligible cost. By 2030, the practical distance between what the standards
deliver and what information theory permits will be the smallest it has ever
been for browser-based real-time communication. The consequence for
application developers is that the infrastructure --- standards, hardware,
and networks --- will no longer be the binding constraint. The remaining
challenge is purely architectural: how to compose these efficient, universal
primitives into applications that fully exploit the capacity and low latency
they provide.

%
% Requires in the preamble: amsmath, amssymb, tabularx, booktabs

\section{Information-Theoretic View}
\label{app:info-theory}

This appendix applies information-theoretic measures to the real-time
web communication stack described in this document. We quantify channel
capacity, source coding efficiency, and effective information rates
across three time horizons: today (March~2026), the near term
(2027--2028), and the medium term (2029--2031). The analysis draws on
the standards, codec parameters, and infrastructure projections
developed in the main text.

\subsection{Channel Capacity of the Transport Layer}

The Shannon-Hartley theorem gives the theoretical maximum information
rate for a channel with additive white Gaussian noise:
\begin{equation}
  C = B \log_2\!\left(1 + \frac{S}{N}\right) \quad \text{[bits/s]}
  \label{eq:shannon}
\end{equation}
where $B$ is the channel bandwidth in Hz and $S/N$ is the
signal-to-noise ratio.

While the physical-layer capacity of a broadband link is governed by
(\ref{eq:shannon}), the \emph{effective} capacity available to the
application is reduced by protocol overhead, congestion control, and
transport inefficiencies. We define the effective application-layer
throughput as:
\begin{equation}
  R_{\text{eff}} = C_{\text{link}} \cdot \eta_{\text{transport}}
    \cdot (1 - \rho_{\text{overhead}})
  \label{eq:reff}
\end{equation}
where $C_{\text{link}}$ is the link-layer capacity,
$\eta_{\text{transport}} \in (0,1]$ is the transport efficiency factor
(capturing congestion control, retransmission, and multiplexing
behavior), and $\rho_{\text{overhead}} \in [0,1)$ is the fractional
protocol overhead (headers, encryption, framing).

\subsubsection{Head-of-Line Blocking and Multiplexed Capacity}

Consider $n$ independent streams multiplexed over a single connection.
Under TCP (HTTP/2), a single lost packet blocks all $n$ streams. The
expected throughput per stream under loss rate $p$ is bounded by:
\begin{equation}
  R_{\text{TCP}}^{(i)} \leq \frac{R_{\text{eff}}}{n}
    \cdot (1 - p)^{n-1}
  \label{eq:hol-tcp}
\end{equation}
because a loss in \emph{any} of the $n$ streams stalls all others. Under
QUIC (RFC~9000), streams are independent at the transport layer, so:
\begin{equation}
  R_{\text{QUIC}}^{(i)} \leq \frac{R_{\text{eff}}}{n}
    \cdot (1 - p)
  \label{eq:hol-quic}
\end{equation}
The ratio of effective per-stream throughput is:
\begin{equation}
  \frac{R_{\text{QUIC}}^{(i)}}{R_{\text{TCP}}^{(i)}}
    = \frac{1}{(1-p)^{n-2}}
  \label{eq:hol-ratio}
\end{equation}
\excursus{Excursus: Why $n = 8$?}{In a typical multi-participant video
conference, each peer contributes multiple concurrent streams: a camera video
track, an audio track, and often a screen-share or data channel. A four-person
call with two media tracks each produces $n = 8$ multiplexed streams over a
single connection --- the reference value used throughout this section. Larger
meetings (e.g.\ 10 participants with simulcast layers) push $n$ higher,
amplifying the HOL blocking effect further.}

For $n = 8$ streams (a typical multi-track video conference) at
$p = 0.01$ (1\% loss), this gives a QUIC advantage factor of
$1/(0.99)^6 \approx 1.062$, a modest 6.2\% improvement. At $p = 0.05$
(5\% loss, common on mobile or congested links), the factor is
$1/(0.95)^6 \approx 1.34$, a 34\% throughput advantage per stream.

\begin{figure}[ht]
\centering
\begin{tikzpicture}[>=Stealth]

% --- Shaded region (drawn FIRST so curves appear on top) ---
\fill[orange!8] (2, 0) rectangle (10, 5.3);
\node[font=\tiny\itshape, text=black!40, anchor=north east] at (9.8, 5.2)
  {mobile / congested};
\node[font=\tiny\itshape, text=black!40, anchor=north west] at (0.2, 5.2)
  {fixed broadband};

% --- Grid lines ---
\foreach \y in {1, 2, 3, 4, 5} {
  \draw[gray!20] (0, \y) -- (10, \y);
}

% --- Axes ---
\draw[->, thick] (0, 0) -- (10.5, 0);
\draw[->, thick] (0, 0) -- (0, 5.5);

% --- Axis titles (positioned to clear tick labels) ---
\node[font=\small] at (5, -0.85) {Packet loss rate $p$};
\node[font=\small, rotate=90] at (-0.9, 2.75) {QUIC advantage factor};

% --- X-axis labels ---
\foreach \x/\lab in {0/0, 2/0.02, 4/0.04, 6/0.06, 8/0.08, 10/0.10} {
  \draw (\x, -0.1) -- (\x, 0.1);
  \node[font=\tiny, anchor=north] at (\x, -0.15) {\lab};
}

% --- Y-axis labels ---
\foreach \y/\lab in {0/1.0, 1/1.1, 2/1.2, 3/1.3, 4/1.4, 5/1.5} {
  \draw (-0.1, \y) -- (0.1, \y);
  \node[font=\tiny, anchor=east] at (-0.15, \y) {\lab};
}

% --- n=4 curve: 1/(1-p)^2 ---
\draw[very thick, blue!70]
  (0, 0) .. controls (2.5, 0.5) and (5, 1.1) ..
  (7.5, 2.0) .. controls (8.5, 2.4) and (9.5, 2.9) ..
  (10, 3.1);
\node[font=\scriptsize, blue!70, anchor=south west] at (9.0, 3.0)
  {$n = 4$};

% --- n=8 curve: 1/(1-p)^6 ---
\draw[very thick, orange!80!black]
  (0, 0) .. controls (1.5, 0.5) and (3, 1.7) ..
  (4, 2.5) .. controls (5, 3.5) and (5.5, 4.0) ..
  (6, 4.7);
\node[font=\scriptsize, orange!80!black, anchor=south] at (5.5, 4.8)
  {$n = 8$};

% --- Reference points ---
\fill[orange!80!black] (1.0, 0.62) circle (2.5pt);
\node[font=\tiny, orange!80!black, anchor=south west] at (1.1, 0.62)
  {6.2\%};

\fill[orange!80!black] (5.0, 3.7) circle (2.5pt);
\node[font=\tiny, orange!80!black, anchor=west] at (5.1, 3.7)
  {34\%};

\fill[blue!70] (1.0, 0.20) circle (2.5pt);

% --- Annotation ---
\node[font=\scriptsize\itshape, text=black!60, align=center,
  anchor=north] at (5.0, -1.0)
  {QUIC recovers per-stream throughput closer to $C/n$ (Shannon fair share)\\
   by eliminating head-of-line blocking (eq.~\ref{eq:hol-ratio}).};

\end{tikzpicture}
\caption{QUIC advantage factor $1/(1-p)^{n-2}$ over TCP for multiplexed
  streams. At 5\% loss with $n=8$ streams, QUIC delivers 34\% more
  per-stream throughput by maintaining independence at the transport layer.}
\label{fig:hol-blocking}
\end{figure}

\subsubsection{Connection Establishment Latency}

Let $\tau$ denote the one-way latency (half the round-trip time). The
information-carrying capacity during connection establishment is zero;
the handshake represents pure overhead. The time to first useful byte
is:
\begin{align}
  T_{\text{TCP+TLS\,1.3}} &= 2\tau_{\text{TCP}} + \tau_{\text{TLS}}
    = 3\tau \quad \text{(1-RTT TLS~1.3)} \label{eq:tcp-setup} \\
  T_{\text{QUIC}} &= \tau \quad \text{(1-RTT)}
    \label{eq:quic-setup} \\
  T_{\text{QUIC,\,0-RTT}} &= 0 \quad \text{(resumption)}
    \label{eq:quic-0rtt}
\end{align}
The QUIC 0-RTT case (\ref{eq:quic-0rtt}) eliminates the handshake
entirely for resumed connections. The information lost during
handshake dead time, relative to steady-state capacity, is:
\begin{equation}
  I_{\text{lost}} = C_{\text{link}} \cdot T_{\text{setup}}
  \label{eq:ilost}
\end{equation}
For a 120~Mbps link with $\tau = 25$~ms (typical fixed broadband in
2026), the TCP+TLS~1.3 handshake wastes $120 \times 10^6 \times 0.075
= 9 \times 10^6$ bits (1.125~MB), while QUIC~1-RTT wastes only
$3 \times 10^6$ bits (375~KB).

\subsection{Source Coding Efficiency}

The rate-distortion function $R(D)$ defines the minimum bit rate
required to represent a source at distortion level $D$. For a
Gaussian source with variance $\sigma^2$ under mean squared error
distortion:
\begin{equation}
  R(D) = \frac{1}{2} \log_2 \frac{\sigma^2}{D}
    \quad \text{[bits/sample]}
  \label{eq:rate-distortion}
\end{equation}

\excursus{Excursus: Why Gaussian source?}{The Gaussian memoryless model is used
because it yields the highest (worst-case) $R(D)$ among all sources of equal
variance. Any codec that approaches the Gaussian bound on real video is
therefore doing at least as well as the bound predicts --- the true $R(D)$ for
spatially and temporally correlated video lies below the Gaussian curve, making
it a conservative benchmark.}

No practical codec achieves $R(D)$;\footnote{The $R(D)$ curve in
Figure~\ref{fig:rate-distortion} is the Gaussian memoryless source bound
(equation~\ref{eq:rate-distortion}), which is an upper bound on the true
$R(D)$ for non-Gaussian sources such as video.  Additionally, $R(D)$ assumes
zero excess-distortion probability, whereas practical codecs permit a
frame-level excess-distortion probability on the order of $10^{-6}$ to
$10^{-8}$.  Both factors allow codec operating points measured on real content
to approach or marginally cross the Gaussian bound.} real codecs operate at
some rate $R_{\text{codec}} > R(D)$. We define the coding efficiency gap as:
\begin{equation}
  \Delta = \frac{R_{\text{codec}} - R(D)}{R(D)}
  \label{eq:coding-gap}
\end{equation}

The codecs in the WebRTC stack can be ordered by their empirical
proximity to $R(D)$ for typical video content at conversational
resolutions (720p, 30~fps). Using published Bj\o ntegaard Delta (BD)
rate measurements:
\begin{equation}
  R_{\text{H.264}} > R_{\text{VP9}} > R_{\text{AV1}}
    > R(D)
  \label{eq:codec-ordering}
\end{equation}

Specifically, the empirical rate savings at equivalent PSNR are
approximately:
\begin{align}
  R_{\text{VP9}}  &\approx 0.65 \cdot R_{\text{H.264}}
    \quad (\sim\!35\% \text{ savings}) \label{eq:vp9-savings} \\
  R_{\text{AV1}}  &\approx 0.70 \cdot R_{\text{VP9}}
    \approx 0.46 \cdot R_{\text{H.264}}
    \quad (\sim\!30\% \text{ over VP9})
  \label{eq:av1-savings}
\end{align}

\begin{figure}[ht]
\centering
\begin{tikzpicture}[>=Stealth]

% --- Axes ---
\draw[->, thick] (0, 0) -- (10.5, 0);
\draw[->, thick] (0, 0) -- (0, 6.5);

% --- Axis titles ---
\node[font=\small] at (5, -0.85) {Distortion $D$};
\node[font=\small, rotate=90] at (-0.9, 3.25) {Rate [Mbps]};

% --- Y-axis rate labels ---
\foreach \y/\lab in {1.0/0.5, 2.0/1.0, 3.0/1.5, 4.0/2.0, 5.0/2.5, 6.0/3.0} {
  \draw (-0.1, \y) -- (0.1, \y);
  \node[font=\tiny, anchor=east] at (-0.15, \y) {\lab};
}

% --- R(D) Shannon bound curve ---
\draw[very thick, green!60!black]
  (1.0, 5.8) .. controls (2.0, 3.5) and (3.5, 2.0) ..
  (5.0, 1.3) .. controls (6.5, 0.8) and (8.0, 0.55) ..
  (9.5, 0.4);
\node[font=\scriptsize, green!60!black, anchor=north west] at (8.5, 0.55)
  {$R(D)$ bound};

% --- Reference distortion line (720p/30fps operating point) ---
\draw[thin, gray, dashed] (4.0, 0) -- (4.0, 6.2);
\node[font=\tiny, gray, anchor=south] at (4.0, 6.2) {720p/30fps};

% --- R(D) value at reference distortion ---
\fill[green!60!black] (4.0, 1.55) circle (2pt);

% --- Codec operating points at reference distortion ---
% H.264: ~2.5 Mbps
\draw[thick, dashed, red!70!black] (0.3, 4.8) -- (9.5, 4.8);
\fill[red!70!black] (4.0, 4.8) circle (3pt);
\node[font=\scriptsize, red!70!black, anchor=west] at (9.6, 4.8) {H.264};

% VP9: ~1.6 Mbps (0.65 * H.264)
\draw[thick, dashed, blue!70] (0.3, 3.3) -- (9.5, 3.3);
\fill[blue!70] (4.0, 3.3) circle (3pt);
\node[font=\scriptsize, blue!70, anchor=west] at (9.6, 3.3) {VP9};

% AV1: ~1.15 Mbps (0.46 * H.264)
\draw[thick, dashed, orange!80!black] (0.3, 2.3) -- (9.5, 2.3);
\fill[orange!80!black] (4.0, 2.3) circle (3pt);
\node[font=\scriptsize, orange!80!black, anchor=west] at (9.6, 2.3) {AV1};

% --- Gap annotations (braces) ---
% H.264 gap
\draw[decorate, decoration={brace, amplitude=4pt, mirror},
  thick, red!50!black]
  (3.6, 1.55) -- (3.6, 4.8)
  node[midway, anchor=east, font=\scriptsize, xshift=-5pt]
  {$\Delta_{\text{H.264}}$};

% AV1 gap
\draw[decorate, decoration={brace, amplitude=4pt},
  thick, orange!60!black]
  (4.4, 1.55) -- (4.4, 2.3)
  node[midway, anchor=west, font=\scriptsize, xshift=5pt]
  {$\Delta_{\text{AV1}}$};

% --- Annotation ---
\node[font=\scriptsize\itshape, text=black!60, align=center,
  anchor=north] at (5.0, -0.5)
  {The Shannon rate-distortion bound $R(D)$ is the theoretical minimum.\\
   AV1 narrows the gap $\Delta$ to ${\sim}0.4$ by 2030 (eq.~\ref{eq:coding-gap}).};

\end{tikzpicture}
\caption{Codec operating rates vs.\ the Shannon rate-distortion bound $R(D)$
  at 720p/30fps. Each codec operates above the theoretical minimum; the gap
  $\Delta$ (eq.~\ref{eq:coding-gap}) shrinks from H.264 through VP9 to AV1,
  approaching but never reaching the bound.}
\label{fig:rate-distortion}
\end{figure}

\subsubsection{Codec Entropy and WebCodecs}

The WebCodecs API (W3C Working Draft) exposes per-frame codec access,
allowing measurement and control of the instantaneous encoding rate.
For a video frame $f_t$ at time $t$, the encoder output has empirical
entropy:
\begin{equation}
  \hat{H}(f_t) = -\sum_{s \in \mathcal{S}}
    p(s \mid f_t) \log_2 p(s \mid f_t)
  \label{eq:frame-entropy}
\end{equation}
where $\mathcal{S}$ is the set of symbols emitted by the entropy coder
(CABAC for H.264/AV1, arithmetic coding for VP9). WebCodecs enables
applications to observe $|E(f_t)|$ (the encoded frame size in bits)
directly, which serves as a practical upper bound on
$\hat{H}(f_t)$:
\begin{equation}
  \hat{H}(f_t) \leq |E(f_t)| \leq R_{\text{codec}} / \text{fps}
  \label{eq:entropy-bound}
\end{equation}

\subsection{Effective Information Rate by Time Horizon}

We now combine transport and source coding measures to estimate the
effective information rate for a representative real-time video
session: a single 720p/30fps video call with one audio channel.

\excursus{Excursus: Why 720p?}{Throughout this appendix, 720p/30fps is used as
a fixed reference for cross-era comparison, not as a prediction that 720p
remains the dominant conferencing resolution. The bandwidth surplus quantified
below is precisely what enables migration to 1080p (feasible by ${\sim}$2028
with AV1 hardware encode at today's 720p bit rates) and 4K (feasible by
${\sim}$2030). The 720p baseline makes the efficiency gains legible: the same
task gets cheaper, freeing capacity for higher quality or more streams.}

Define the \emph{net information rate} as:
\begin{equation}
  I_{\text{net}} = R_{\text{codec}} \cdot \eta_{\text{transport}}
    \cdot (1 - \rho_{\text{overhead}})
    \cdot (1 - p_{\text{loss}})
  \label{eq:inet}
\end{equation}

\begin{table}[ht]
\centering
\small
\renewcommand{\arraystretch}{1.4}
\begin{tabularx}{\textwidth}{lXXX}
\toprule
\textbf{Parameter} & \textbf{March 2026} & \textbf{2027--2028} &
  \textbf{2029--2031} \\
\midrule
Median link capacity &
  120~Mbps &
  $\sim$150~Mbps &
  200--300~Mbps \\
Dominant codec (WebRTC) &
  VP9 / H.264 &
  VP9 / AV1 (HW) &
  AV1 (HW default) \\
Codec rate (720p/30fps) &
  1.5--2.5~Mbps &
  1.0--1.8~Mbps &
  0.7--1.2~Mbps \\
$\eta_{\text{transport}}$ &
  0.92 (QUIC) &
  0.93 &
  0.94 \\
$\rho_{\text{overhead}}$ (SRTP+QUIC) &
  0.06 &
  0.05 &
  0.04 \\
Typical loss $p$ &
  0.01--0.03 &
  0.01--0.02 &
  0.005--0.01 \\
$I_{\text{net}}$ (720p video) &
  1.3--2.2~Mbps &
  0.9--1.6~Mbps &
  0.6--1.1~Mbps \\
$I_{\text{net}}$ as \% of $C_{\text{link}}$ &
  1.1--1.8\% &
  0.6--1.1\% &
  0.2--0.6\% \\
\bottomrule
\end{tabularx}
\caption{Effective information rate for a 720p/30fps WebRTC video stream.
  The declining ratio $I_{\text{net}} / C_{\text{link}}$ reflects both
  improved codec efficiency and growing link capacity, freeing bandwidth
  for higher resolutions, additional streams, or concurrent data
  channels.}
\label{tab:inet}
\end{table}

The declining ratio $I_{\text{net}} / C_{\text{link}}$ in
Table~\ref{tab:inet} is the central quantitative observation: by 2030, a
720p video call consumes less than 0.5\% of available link capacity. The
surplus capacity enables qualitative changes --- 4K video conferencing,
multi-stream layouts, or concurrent data channels via WebTransport ---
that are quantitatively infeasible in 2026.

\begin{figure}[ht]
\centering
\begin{tikzpicture}[>=Stealth]

% --- Axes ---
\draw[->, thick] (0, 0) -- (11, 0);
\draw[->, thick] (0, 0) -- (0, 5.5);

% --- Axis title ---
\node[font=\small, rotate=90] at (-0.9, 2.75) {Mbps};

% --- Y-axis labels ---
\foreach \y/\lab in {0/0, 1/50, 2/100, 3/150, 4/200, 5/250} {
  \draw (-0.1, \y) -- (0.1, \y);
  \node[font=\tiny, anchor=east] at (-0.15, \y) {\lab};
}

% --- Group: March 2026 ---
% C_link = 120 Mbps -> y = 2.4
\fill[blue!12] (0.8, 0) rectangle (2.4, 2.4);
\draw[blue!50, thick] (0.8, 0) rectangle (2.4, 2.4);
% I_net = 1.75 Mbps (midpoint) -> y = 0.035
\fill[orange!60] (0.8, 0) rectangle (2.4, 0.07);
\draw[orange!80!black, thick] (0.8, 0) rectangle (2.4, 0.07);
\node[font=\footnotesize, anchor=south] at (1.6, 2.5) {120};
\node[font=\tiny, orange!80!black] at (1.6, -0.4) {1.5\%};
\node[font=\small, anchor=north] at (1.6, -0.7) {2026};

% --- Group: 2027--2028 ---
% C_link = 150 Mbps -> y = 3.0
\fill[blue!12] (4.0, 0) rectangle (5.6, 3.0);
\draw[blue!50, thick] (4.0, 0) rectangle (5.6, 3.0);
% I_net = 1.25 Mbps (midpoint) -> y = 0.025
\fill[orange!60] (4.0, 0) rectangle (5.6, 0.05);
\draw[orange!80!black, thick] (4.0, 0) rectangle (5.6, 0.05);
\node[font=\footnotesize, anchor=south] at (4.8, 3.1) {150};
\node[font=\tiny, orange!80!black] at (4.8, -0.4) {0.8\%};
\node[font=\small, anchor=north] at (4.8, -0.7) {2027--28};

% --- Group: 2029--2031 ---
% C_link = 250 Mbps -> y = 5.0
\fill[blue!12] (7.2, 0) rectangle (8.8, 5.0);
\draw[blue!50, thick] (7.2, 0) rectangle (8.8, 5.0);
% I_net = 0.85 Mbps (midpoint) -> y = 0.017
\fill[orange!60] (7.2, 0) rectangle (8.8, 0.034);
\draw[orange!80!black, thick] (7.2, 0) rectangle (8.8, 0.034);
\node[font=\footnotesize, anchor=south] at (8.0, 5.1) {250};
\node[font=\tiny, orange!80!black] at (8.0, -0.4) {0.3\%};
\node[font=\small, anchor=north] at (8.0, -0.7) {2029--31};

% --- Legend ---
\node[anchor=west, font=\footnotesize] at (0.5, 5.8) {%
  \tikz[baseline=-0.5ex]{\fill[blue!12] (0,0) rectangle (0.3,0.3);
    \draw[blue!50, thick] (0,0) rectangle (0.3,0.3);}
  ~$C_{\text{link}}$ (channel capacity) \qquad
  \tikz[baseline=-0.5ex]{\fill[orange!60] (0,0) rectangle (0.3,0.3);
    \draw[orange!80!black, thick] (0,0) rectangle (0.3,0.3);}
  ~$I_{\text{net}}$ (720p stream)};

% --- Annotation ---
\node[font=\scriptsize\itshape, text=black!60, align=center,
  anchor=north] at (5.5, -1.2)
  {A single 720p call uses a vanishing fraction of Shannon-derived\\
   channel capacity, freeing bandwidth for 4K, multi-stream, and data channels.};

\end{tikzpicture}
\caption{Channel capacity $C_{\text{link}}$ vs.\ effective information rate
  $I_{\text{net}}$ for a single 720p/30fps stream across time horizons. The
  orange bar (stream rate) is barely visible against the blue bar (capacity),
  illustrating the growing surplus as codecs improve and bandwidth increases.}
\label{fig:inet-capacity}
\end{figure}

\subsection{Security Entropy and Key Space}

The DTLS-SRTP security architecture (RFCs~8827, 3711) protects all
WebRTC media. Information-theoretic security requires that the key
entropy exceed the information content of the protected material.

For a DTLS~1.3 handshake with ECDHE on P-256, the key agreement
provides:
\begin{equation}
  H_{\text{key}} = 128 \text{~bits of security}
  \label{eq:key-entropy}
\end{equation}

The SRTP master key is 128 bits (AES-128-GCM), giving an exhaustive
search cost of $2^{128}$ operations. The total information content of a
one-hour 720p video call at 2~Mbps is:
\begin{equation}
  I_{\text{call}} = 2 \times 10^6 \times 3600
    = 7.2 \times 10^9 \text{~bits}
    \approx 2^{32.7} \text{~bits}
  \label{eq:call-content}
\end{equation}

The ratio of key entropy to protected content:
\begin{equation}
  \frac{H_{\text{key}}}{I_{\text{call}}}
    = \frac{128}{32.7} \approx 3.91
  \label{eq:security-margin}
\end{equation}

This ratio remains comfortably above~1 even for much longer sessions.
For a continuous 24-hour 4K AV1 stream at 15~Mbps:
\begin{equation}
  I_{\text{stream}} = 15 \times 10^6 \times 86400
    = 1.296 \times 10^{12} \text{~bits}
    \approx 2^{40.2} \text{~bits}
\end{equation}
yielding $H_{\text{key}} / I_{\text{stream}} \approx 3.18$, still
well above the threshold.

\subsubsection{End-to-End Encryption via Encoded Transforms}

The RTCRtpScriptTransform API (Encoded Transforms) enables
application-layer E2EE within the WebRTC pipeline. From an
information-theoretic perspective, this introduces a second encryption
layer with its own key, creating a \emph{nested channel}:
\begin{equation}
  X \xrightarrow{\;K_{\text{E2EE}}\;}
  E_1(X) \xrightarrow{\;K_{\text{SRTP}}\;}
  E_2(E_1(X))
  \label{eq:nested}
\end{equation}

The mutual information between the ciphertext and plaintext, given
neither key, satisfies:
\begin{equation}
  I\!\left(X;\, E_2(E_1(X))\right) \leq
    \min\!\big(H(K_{\text{E2EE}})^{-1},\;
    H(K_{\text{SRTP}})^{-1}\big)
    \cdot I_{\text{call}}
  \label{eq:mutual-e2ee}
\end{equation}

In practice, with $H(K_{\text{E2EE}}) = H(K_{\text{SRTP}}) = 128$
bits, this quantity is negligible ($\ll 1$ bit for any practical
session length), confirming that the nested encryption provides
information-theoretic confidentiality well beyond computational
security guarantees.

Today (March~2026), Encoded Transforms remain Chrome-only. By
2027--2028 (Section~8.3), cross-browser standardization will make
this nested channel universally available. By 2029--2031, SFrame
for group calls will extend this to $n$-party sessions, where the
mutual information between any non-member and the group plaintext
remains negligible:
\begin{equation}
  I\!\left(X;\, E(X) \;\middle|\; K_i \notin
    \{K_1, \ldots, K_n\}\right) \approx 0
  \label{eq:sframe}
\end{equation}

\subsection{WebTransport: Unreliable Datagrams and Channel Capacity}

WebTransport (RFC~9297) introduces unreliable datagrams to browser-server
communication. For latency-sensitive applications (game state, sensor
telemetry), the relevant measure is not throughput but
\emph{freshness-weighted information rate}. Define the value of a
datagram as exponentially decaying with delivery latency $\delta$:
\begin{equation}
  V(\delta) = e^{-\lambda \delta}
  \label{eq:freshness}
\end{equation}
where $\lambda > 0$ is an application-specific decay rate. The
effective information rate under unreliable delivery is:
\begin{equation}
  I_{\text{WT}} = R_{\text{send}} \cdot (1 - p_{\text{loss}})
    \cdot \mathbb{E}\!\left[e^{-\lambda \delta}\right]
  \label{eq:iwt}
\end{equation}

Under reliable delivery (TCP/WebSocket), a retransmission delays
subsequent data, so:
\begin{equation}
  I_{\text{WS}} = R_{\text{send}}
    \cdot \mathbb{E}\!\left[e^{-\lambda(\delta + p \cdot 2\tau)}\right]
  \label{eq:iws}
\end{equation}

The ratio $I_{\text{WT}} / I_{\text{WS}}$ exceeds 1 when:
\begin{equation}
  (1 - p) \cdot e^{\lambda \cdot p \cdot 2\tau} > 1
  \quad \Longleftrightarrow \quad
  \lambda > \frac{-\ln(1 - p)}{2 p \tau}
  \label{eq:wt-advantage}
\end{equation}

For $p = 0.02$ and $\tau = 25$~ms, this gives $\lambda > 20.2$~s$^{-1}$
(half-life $< 34$~ms).

\excursus{Excursus: Why $p = 0.02$ and $\tau = 25$\,ms?}{These values represent
median fixed-broadband conditions in 2026 (Table~\ref{tab:inet}). On mobile or
congested links ($p \approx 0.05$, $\tau \approx 50$\,ms), the threshold
$\lambda^*$ drops to ${\sim}10$\,s$^{-1}$, broadening WebTransport's advantage
to applications with data half-lives up to ${\sim}70$\,ms. The fixed-broadband
reference is thus conservative: mobile scenarios favor unreliable delivery even
more strongly.}

Any application where data older than
$\sim$30~ms is useless benefits from WebTransport's unreliable
mode over WebSocket --- precisely the regime of interactive gaming, live
cursor tracking, and real-time sensor feeds.

As of March~2026, Safari does not support WebTransport
(Section~5), limiting coverage to $\sim$82\%. By 2027--2028
(Section~8.2), universal support will allow applications to rely on
(\ref{eq:iwt}) without fallback paths.

\begin{figure}[ht]
\centering
\begin{tikzpicture}[>=Stealth]

% --- Axes ---
\draw[->, thick] (0, 0) -- (10.5, 0);
\draw[->, thick] (0, 0) -- (0, 5.5);

% --- Axis titles (positioned to clear tick labels) ---
\node[font=\small] at (5, -0.75)
  {Freshness decay rate $\lambda$ [s$^{-1}$]};
\node[font=\small, rotate=90] at (-1.1, 2.75)
  {$I_{\text{WT}} / I_{\text{WS}}$};

% --- X-axis labels ---
\foreach \x/\lab in {0/0, 1.67/10, 3.33/20, 5.0/30, 6.67/40, 8.33/50, 10/60} {
  \draw (\x, -0.1) -- (\x, 0.1);
  \node[font=\tiny, anchor=north] at (\x, -0.15) {\lab};
}

% --- Y-axis labels (0.96 to 1.06, each 0.02 = 1 unit) ---
\foreach \y/\lab in {0/0.96, 1/0.98, 2/1.00, 3/1.02, 4/1.04, 5/1.06} {
  \draw (-0.1, \y) -- (0.1, \y);
  \node[font=\tiny, anchor=east] at (-0.15, \y) {\lab};
}

% --- Ratio = 1.0 line (y=2 maps to ratio 1.0) ---
\draw[thin, gray, dashed] (0, 2) -- (10, 2);

% --- Shade WebTransport advantage region ---
% lambda* = 20.2 -> x = 3.37
\fill[green!10] (3.37, 2) rectangle (10, 5.3);

% --- I_WT/I_WS curve: 0.98*exp(lambda*0.001) ---
% Y range: 0.96-1.06 (each 0.02 = 1 unit). lambda=0: y=1, lambda*=20.2: y=2, lambda=60: y=4.05
\draw[very thick, orange!80!black]
  (0, 1.0) .. controls (1.5, 1.3) and (2.5, 1.7) ..
  (3.37, 2.0) .. controls (4.5, 2.4) and (6, 2.8) ..
  (8, 3.4) .. controls (9, 3.7) and (9.5, 3.85) ..
  (10, 4.05);

% --- Threshold line ---
\draw[thick, dashed, red!60!black] (3.37, 0) -- (3.37, 5.2);
\node[font=\scriptsize, red!60!black, anchor=south, rotate=90] at (3.17, 3.5)
  {$\lambda^* = 20.2$ s$^{-1}$};

% --- Application regime labels ---
\node[font=\scriptsize\itshape, text=blue!60, anchor=north] at (1.5, 5.3)
  {video conf};
\node[font=\scriptsize\itshape, text=green!50!black, anchor=north] at (7.0, 5.3)
  {gaming / cursor / sensors};

% --- Mark advantage/disadvantage ---
\node[font=\tiny, text=red!50!black, anchor=north east] at (3.2, 1.8)
  {WebSocket wins};
\node[font=\tiny, text=green!50!black, anchor=north west] at (3.5, 4.5)
  {WebTransport wins};

% --- Annotation ---
\node[font=\scriptsize\itshape, text=black!60, align=center,
  anchor=north] at (5.0, -1.0)
  {For $p = 0.02$, $\tau = 25$~ms. Applications where data older than\\
   ${\sim}$30~ms is stale gain from unreliable delivery (eq.~\ref{eq:wt-advantage}).};

\end{tikzpicture}
\caption{Freshness-weighted information rate ratio $I_{\text{WT}} / I_{\text{WS}}$
  as a function of decay rate $\lambda$. Above the threshold $\lambda^*$,
  WebTransport's unreliable datagrams deliver more useful information per
  unit time than WebSocket's reliable delivery, because retransmission
  wastes capacity on stale data.}
\label{fig:freshness}
\end{figure}

\subsection{Projections Summary and Practical Implications}

\begin{table}[ht]
\centering
\small
\renewcommand{\arraystretch}{1.4}
\begin{tabularx}{\textwidth}{lXXX}
\toprule
\textbf{Measure} & \textbf{March 2026} & \textbf{2027--2028} &
  \textbf{2029--2031} \\
\midrule
Channel capacity $C_{\text{link}}$ &
  120~Mbps &
  150~Mbps &
  250~Mbps \\
HOL blocking factor $(1\!-\!p)^{n-2}$ at $n\!=\!8$ &
  0.94 ($p\!=\!0.01$) &
  0.94 &
  0.97 ($p\!=\!0.005$) \\
Source coding gap $\Delta$ &
  VP9: $\sim$0.8 &
  AV1 HW: $\sim$0.5 &
  AV1 opt: $\sim$0.4 \\
$I_{\text{net}}$ per 720p stream &
  1.3--2.2~Mbps &
  0.9--1.6~Mbps &
  0.6--1.1~Mbps \\
Max 720p streams in $C_{\text{link}}$ &
  $\sim$55--90 &
  $\sim$95--165 &
  $\sim$230--415 \\
E2EE availability &
  Chrome only &
  Cross-browser &
  SFrame $n$-party \\
WebTransport coverage &
  82\% &
  $\sim$100\% &
  100\% \\
Freshness threshold $\lambda^*$ &
  20.2~s$^{-1}$ &
  $\sim$18~s$^{-1}$ &
  $\sim$15~s$^{-1}$ \\
\bottomrule
\end{tabularx}
\caption{Information-theoretic measures across time horizons. $\Delta$
  is the coding efficiency gap (\ref{eq:coding-gap}), $I_{\text{net}}$
  is the net information rate (\ref{eq:inet}), and $\lambda^*$ is the
  freshness decay threshold above which WebTransport dominates WebSocket
  (\ref{eq:wt-advantage}).}
\label{tab:projections}
\end{table}

The trajectory is clear: the combination of improving source coding
(AV1 hardware), expanding channel capacity (broadband growth), and
maturing transport protocols (QUIC universal, WebTransport universal)
drives the \emph{information cost per unit of communicated experience}
steadily downward. The remaining question is not whether the
Shannon capacity suffices --- it already does, by orders of magnitude
--- but how efficiently the protocol stack converts that capacity into
delivered, fresh, secure information at the application layer. By
2030, the gap between theoretical capacity and realized throughput will
be the narrowest it has ever been for browser-based real-time
communication.

\subsubsection{What the Theory Means for Applications}

The equations and figures in this appendix are not purely academic. Each
theoretical result maps to a concrete design decision that application
developers face today and over the next five years.

\textbf{Codec selection determines bandwidth cost, not quality.}
The rate-distortion analysis (Section~A.2, Figure~\ref{fig:rate-distortion})
shows that AV1 operates at roughly 46\% of H.264's bit rate for equivalent
perceptual quality at 720p. In practice, this means a WebRTC application
switching from H.264 to AV1 can either halve its bandwidth consumption at
the same quality or substantially increase resolution (e.g., 720p to 1080p)
within the same bandwidth envelope. By 2030, with AV1 hardware encoding
ubiquitous and the coding gap $\Delta$ approaching 0.4, further codec
improvements yield diminishing returns --- the stack will be operating close
enough to the Shannon rate-distortion bound that the next generation of
efficiency gains will come from transport and architecture, not codecs alone.

\textbf{QUIC's advantage is situational but decisive on bad networks.}
The head-of-line blocking analysis (Section~A.1.1,
Figure~\ref{fig:hol-blocking}) shows that QUIC's per-stream throughput
advantage over TCP is modest on good networks (6\% at 1\% loss) but
substantial on degraded ones (34\% at 5\% loss). The practical implication
is that applications targeting mobile users, emerging markets, or congested
enterprise networks --- precisely the environments where real-time
communication is hardest --- benefit disproportionately from QUIC-based
transports. For a video conference with 8 participants on a cellular
connection experiencing 3--5\% packet loss, the difference between QUIC and
TCP is the difference between usable and degraded video.

\textbf{Connection latency is a solved problem for returning users.}
The handshake analysis (Section~A.1.2) shows that QUIC 0-RTT eliminates
connection establishment overhead entirely for resumed connections. For a
WebTransport application that a user opens daily --- a collaboration tool,
a live dashboard, a recurring meeting client --- the time to first useful
byte drops to zero. The 1.125~MB of wasted capacity during a TCP+TLS
handshake on a 120~Mbps link may seem small, but for latency-sensitive
applications where the first 100~ms of perceived responsiveness determines
user experience, the difference is measurable.

\textbf{Bandwidth surplus enables architectural ambition.}
The effective information rate analysis (Section~A.3,
Figure~\ref{fig:inet-capacity}) reveals the most consequential practical
finding: a single 720p video stream consumes a vanishing fraction of
available link capacity (under 0.5\% by 2030). This surplus is not merely
``headroom'' --- it is the resource that makes new application architectures
feasible. A 250~Mbps link can theoretically carry 230--415 simultaneous 720p
streams. In practice, this means applications can afford to run multiple
concurrent media flows (video, screen share, spatial audio), data channels
(collaborative editing, file transfer), and telemetry streams without
contending for bandwidth. The constraint shifts from ``can the network carry
this?''\ to ``can the application manage this complexity?''

\textbf{End-to-end encryption adds negligible overhead.}
The security analysis (Section~A.4) confirms that the nested DTLS-SRTP plus
E2EE architecture imposes no meaningful information-theoretic cost. The key
entropy ($2^{128}$) exceeds the information content of any practical session
by a factor of at least 3. For application developers, this means E2EE can
be treated as a default rather than a premium feature --- there is no
bandwidth, latency, or capacity penalty for encrypting at the application
layer on top of the transport-layer encryption that WebRTC already mandates.
The only constraint is API availability, which moves from Chrome-only (2026)
to universal (2028).

\textbf{WebTransport unlocks real-time data that WebSocket cannot serve.}
The freshness analysis (Section~A.5, Figure~\ref{fig:freshness}) quantifies
when unreliable delivery outperforms reliable delivery: any data with a
useful half-life below approximately 34~ms benefits from WebTransport's
unreliable datagrams. This covers a wide class of interactive applications
--- game state synchronization, live cursor and pointer sharing, real-time
sensor feeds, collaborative whiteboard strokes --- where retransmitting a
stale update is worse than dropping it. Before WebTransport, these
applications had to choose between WebSocket (reliable but stale under loss)
and WebRTC DataChannel (unreliable mode available but requiring the full
peer-to-peer ICE/DTLS machinery). WebTransport provides the unreliable
mode with a simple client-server model over HTTP/3, removing the
architectural overhead.

\textbf{The overall picture: the stack converges on efficiency.}
Taken together, the information-theoretic measures trace a consistent
trajectory. The protocol stack is converging toward the theoretical limits
established by Shannon: codecs approach the rate-distortion bound, transports
eliminate unnecessary blocking and handshake overhead, and encryption adds
negligible cost. By 2030, the practical distance between what the standards
deliver and what information theory permits will be the smallest it has ever
been for browser-based real-time communication. The consequence for
application developers is that the infrastructure --- standards, hardware,
and networks --- will no longer be the binding constraint. The remaining
challenge is purely architectural: how to compose these efficient, universal
primitives into applications that fully exploit the capacity and low latency
they provide.
